\chapter{Cellular Generations: From AMPS to 5G}
\label{ch:cellular}

\begin{chapterintro}
In 1980 a mobile telephone was a forty-pound box in the boot of a car, and a large city
could serve a few thousand of them. By 2025 there were more mobile subscriptions than people
on Earth, and the device in a teenager's pocket moved more data in an evening than the whole
mobile network of 1990 carried in a year. This chapter tells how that happened, one
generation at a time, and takes each generation apart to show its engineering. It starts with
the pre-cellular radiotelephone and the cellular idea, then examines the analog AMPS family
and its security failures, the GSM system in depth (frame structure, bursts, logical channels,
speech coding, frequency hopping, the SIM and the A5 ciphers), its 2G rivals including
Qualcomm's CDMA, and the packet-data retrofits GPRS and EDGE. The 3G section covers the
politics of IMT-2000, WCDMA's fast power control and soft handover, CDMA2000 and HSPA. The
4G section works through LTE's flat all-IP architecture, its frame and reference signals, its
channels, HARQ, link adaptation and MIMO modes, LTE-Advanced and VoLTE. The 5G section covers
the IMT-2020 requirements, NR's flexible numerology and beam-based access, LDPC and polar codes,
massive MIMO in the field, non-standalone and standalone deployment, the service-based core,
network slicing, the disaggregated RAN and the releases from 15 to 18. The chapter closes with
a side-by-side comparison and the lessons, technical and economic, of forty years of cellular
engineering. The goal throughout is not to list features but to explain \emph{why} each
generation was designed as it was.
\end{chapterintro}

%======================================================================
\section{Four decades, five generations}
\label{sec:ch21:overview}
%======================================================================

Every engineer who works on cellular systems eventually notices that the standards read like
geology. Each generation is a stratum laid down over the previous one, and each carries the
fossils of the problems its designers were trying to solve. GSM's 4.615\,ms frame exists
because a 13\,MHz crystal divided by 48 gives 270.833\,kb/s and eight users were judged the
right number to share a 200\,kHz carrier. LTE's 15\,kHz subcarrier spacing exists because it
had to tolerate a few hundred hertz of Doppler at 2\,GHz and a few microseconds of delay
spread at a sensible cyclic-prefix overhead. NR's 30\,kHz spacing at 3.5\,GHz exists because
the phase noise of an affordable oscillator grows with carrier frequency. Learn the reasons
and the parameters stop being arbitrary numbers to memorise.

\subsection{What a ``generation'' is}

The word \term{generation} is partly engineering and partly marketing. In engineering terms a
new generation is a new air interface that cannot be received by the previous generation's
phones: analog FM (1G), digital TDMA or narrowband CDMA (2G), wideband CDMA (3G), OFDMA with
an all-IP core (4G), and flexible-numerology OFDM with a cloud-native core (5G). In
regulatory terms the generations from 3G onward are defined by the
\term{International Telecommunication Union} (ITU) through the IMT family of recommendations:
IMT-2000 (3G), IMT-Advanced (4G) and IMT-2020 (5G), each with a list of minimum requirements
that a candidate technology must meet. In marketing terms a generation is whatever operators
are prepared to put next to the signal-strength bars, which is why the industry has seen
``3.5G'', ``3.9G'', ``4G LTE'' (for systems that did not meet the IMT-Advanced requirements)
and ``5G E'' (for LTE-Advanced).

The cadence has been remarkably steady: a new generation about every ten years, with the first
commercial networks in 1979--83 (1G), 1991--95 (2G), 2001--03 (3G), 2009--10 (4G) and 2019
(5G), and the next expected around 2030. The steadiness is not a coincidence. A generation
takes about five years to research and standardise, and an operator needs about ten years to
deploy it widely and recover its investment. Silicon improves enough in a decade to make
affordable what was impossible ten years before: a GSM handset's equaliser would have filled a
rack in 1980, and an LTE phone's turbo decoder would have dissipated more than its battery
could supply in 2000.

\mdcfig{ch21_growth}{(a) Global mobile subscriptions by the most advanced technology they
use. 2G dominated for two decades and peaked around 2012; 4G became the majority technology
around 2018; 5G passed a billion subscriptions in 2022. (b) Global mobile data traffic, which
grew by roughly 60\% a year from 2010 to 2024, a factor of about 600 in fourteen years.
Both panels are approximate, compiled from ITU statistics and Ericsson Mobility Reports.}

Figure~\ref{fig:ch21_growth} shows the result. Subscriptions grew from about ten million in
1990 to about 8.8~billion in 2025, and traffic grew from a negligible amount of circuit-switched
voice to well over a hundred exabytes a month. The two curves tell different stories.
Subscription growth was a 1990s and 2000s phenomenon driven by cheap 2G phones and prepaid
billing; data growth is a 2010s phenomenon driven by smartphones and video. Each generation was
designed for the problem of its decade: 1G for mobility itself, 2G for capacity, security and
roaming in voice, 3G for data that turned out to arrive later than planned, 4G for the data
that finally arrived, and 5G for both more data and new kinds of user.

\begin{keyidea}[Each generation fixes its predecessor's bottleneck]
The history of cellular is a sequence of bottlenecks removed. 1G proved cellular reuse but ran
out of capacity and was trivially cloned; 2G went digital for capacity and security. 2G was
circuit-switched voice; 3G built a packet core and a wideband air interface. 3G's architecture
put scheduling and retransmission far from the radio, so its data was slow to react; HSPA and
then LTE moved them into the base station and flattened the network. LTE served smartphones
well but had a single numerology, always-on reference signals and a monolithic core; 5G made
the numerology scalable, the signalling lean and the core a set of software services. When you
study any generation, ask what problem in the previous one it was built to solve.
\end{keyidea}

\subsection{The engineering problems every generation solves}

All five generations solve the same list of problems, with different tools. Keep this list in
mind as a template for the rest of the chapter.
\begin{description}[style=nextline,leftmargin=1.2em,font=\sffamily\bfseries\small\color{navy}]
\item[Sharing the spectrum.] Among users in one cell (FDMA, TDMA, CDMA, OFDMA) and among cells
(frequency reuse, power control, interference coordination). Chapter~\ref{ch:multipleaccess}
develops the theory; this chapter shows the choices.
\item[Beating the channel.] Multipath, fading and Doppler (Chapter~\ref{ch:channels}) are
handled with equalisers (GSM), rake receivers (CDMA), cyclic-prefix OFDM (LTE, NR), coding and
interleaving, diversity and, increasingly, multiple antennas.
\item[Finding and keeping the network.] A phone must find a cell, synchronise to it, register,
be paged, and hand over between cells at speed (Chapter~\ref{ch:sync} covers cell search).
\item[Mobility management.] The network must know roughly where each phone is (location or
tracking areas) and route calls and packets to it as it moves, including across operators and
countries (roaming).
\item[Security.] Subscribers must be authenticated, traffic kept confidential, and identities
protected from tracking; each generation's security reflects the threats of its day and has
been broken in instructive ways.
\item[Economy.] A generation must be cheap enough to build, operate and buy. Battery life,
handset cost, base-station power and spectrum price constrain every design decision.
\end{description}

%======================================================================
\section{Before cellular: the radiotelephone and the cellular idea}
\label{sec:ch21:zerog}
%======================================================================

\subsection{MTS and IMTS}

The first public mobile telephone service in the United States, the Bell System's
\term{Mobile Telephone Service} (MTS), opened in St.~Louis in June 1946. It used a single
high-power transmitter on a tall building to cover a whole metropolitan area, a handful of
FM channels around 150\,MHz, and an operator who connected each call by hand. The equipment
filled much of a car boot and the user pressed a button to talk, as on a two-way radio. The
\term{Improved Mobile Telephone Service} (IMTS) of the mid-1960s added full-duplex operation,
direct dialling and automatic channel selection, but kept the architecture: one big
transmitter per city and a few dozen channels at most. Because every channel was used once
across the entire city, capacity was fixed by the number of channels, not by the number of
users. In the 1970s the New York IMTS system reportedly had a few dozen channels serving
around two thousand subscribers, with thousands more on a waiting list, and a caller might
wait many minutes for a free channel. The other 0G systems around the world (the Nordic MTA
and MTD, Germany's A-Netz and B-Netz) had the same structure and the same ceiling.

\subsection{The cellular concept}

The solution was proposed inside Bell Labs as early as 1947, in an internal memorandum by
Douglas H.~Ring, and developed through the 1960s and 1970s by Richard Frenkiel, Joel Engel and
colleagues. Replace the single high transmitter with many low-power base stations, each
covering a small \term{cell}. Divide the channels into groups, give each cell one group, and
reuse the same group in cells far enough apart that their interference is tolerable. When a
mobile moves from one cell to the next, switch its call to a channel of the new cell, a
\term{handoff} (the American term) or \term{handover} (the European one). Capacity now scales
with the number of cells: halve the cell radius and the number of cells in the city
quadruples, and so does the number of simultaneous calls.

Three ideas make this work, and every later generation inherits them.
\begin{enumerate}
\item \textbf{Frequency reuse} with a \term{cluster size} $N$: for hexagonal cells of radius $R$
the co-channel reuse distance is $D=R\sqrt{3N}$, and with path-loss exponent $n$ the
signal-to-interference ratio at the cell edge is roughly $(3N)^{n/2}/6$. The derivation, the
geometry and the trunking theory of Erlang that sizes each cell are developed in
Chapter~\ref{ch:multipleaccess}; Lab~11 simulates them.
\item \textbf{Cell splitting}: when traffic grows, add base stations and shrink the cells.
\item \textbf{Handoff}: a mobile switching centre that tracks signal quality and moves calls
between cells without the user noticing.
\end{enumerate}
The cost was complexity: many base stations, a switching centre able to reroute calls in
mid-conversation, and mobiles able to retune on command. The technology to build all this
cheaply---the microprocessor and the frequency synthesiser---arrived in the 1970s, and the
spectrum arrived when the US Federal Communications Commission (FCC), after more than a decade
of proceedings, allocated 40\,MHz near 850\,MHz to cellular service in 1981 (raised to 50\,MHz
in 1986).

\begin{history}[The call that started a race]
On 3~April 1973, Martin Cooper of Motorola stood on Sixth Avenue in Manhattan with a
prototype handheld telephone, the DynaTAC, and placed a call through a base station Motorola
had set up on a nearby building. He called Joel Engel, who led the competing cellular effort
at Bell Labs, to tell him that he was speaking from a portable cellular phone. The prototype
weighed about a kilogram and gave around half an hour of talk time after a recharge of about
ten hours. The point of the stunt was regulatory as much as technical: AT\&T's plan was for
car-mounted mobile telephones, and Motorola wanted the FCC to see that the future was a phone
carried by a person. A decade later, in 1983, the commercial DynaTAC 8000X went on sale for
about \$3,995. The two visions---the car phone and the personal handset---are both visible in
the AMPS standard that followed: its power classes run from several watts for vehicle
installations down to 0.6\,W for handhelds.
\end{history}

%======================================================================
\section{First generation: AMPS and its cousins}
\label{sec:ch21:1g}
%======================================================================

\subsection{AMPS}

The \term{Advanced Mobile Phone System} (AMPS), developed at Bell Labs and tested in Chicago
from 1978, opened commercially in Chicago in October 1983. It was an analog FM system in the
825--845\,MHz (mobile transmit) and 870--890\,MHz (base transmit) bands, later extended to
824--849 and 869--894\,MHz. Its main parameters, which defined North American cellular for
twenty years, are as follows.
\begin{itemize}
\item \textbf{Channels.} 30\,kHz channel spacing and a fixed 45\,MHz duplex separation between
the mobile and base frequencies: \term{frequency-division duplex} (FDD). The band was split
between two licensees per market, the ``A'' (non-wireline) and ``B'' (wireline) carriers, each
with 416 channel pairs in the extended allocation, of which 21 were control channels.
\item \textbf{Voice modulation.} FM with a peak deviation of 12\,kHz, preceded by a 2:1
syllabic compandor and pre-emphasis (Chapter~\ref{ch:analog}). Carson's rule with a 3\,kHz
audio band gives about 30\,kHz, the channel spacing.
\item \textbf{Supervisory audio tone.} Each base station added one of three
\term{supervisory audio tones} (SAT) at 5970, 6000 or 6030\,Hz, above the voice band, and the
mobile transponded it back. A mobile that heard the wrong SAT knew it was receiving a
co-channel interferer rather than its own base station, and a base that lost the SAT knew the
link had failed. The SAT is in effect a two-bit colour code for co-channel cells.
\item \textbf{Signalling tone.} A 10\,kHz \term{signalling tone} (ST), sent by the mobile,
signalled on-hook, flash and handoff confirmation.
\item \textbf{Control channels.} Dedicated control channels carried digital messages at
10\,kb/s using Manchester-coded binary FSK with $\pm8$\,kHz deviation, protected by BCH codes
((40,28) on the forward link, (48,36) on the reverse) with each word repeated five times and
decoded by majority vote. During a call, short digital messages, such as a handoff order,
were sent on the voice channel by \term{blank-and-burst}: the audio was muted for about
a tenth of a second while the data burst went through, which users heard as a click.
\item \textbf{Reuse.} Seven-cell clusters with omnidirectional antennas, or more commonly
seven-cell clusters with three 120\textdegree{} sectors each (21 channel groups), chosen
because FM voice needs roughly 17--18\,dB of signal-to-interference ratio for acceptable
quality.
\end{itemize}
Handoff was network-controlled. Each cell site had a \term{locating receiver} that scanned the
channels used in neighbouring cells and measured the mobile's signal strength; when the
serving cell's signal fell below a threshold, the \term{mobile telephone switching office}
(MTSO) chose the neighbour with the strongest measurement and ordered the mobile to retune.

\begin{worked}[How many subscribers can an AMPS cell serve?]
One AMPS licensee has 416 channel pairs, of which 21 carry control, leaving 395 voice channels.
With a 7-cell cluster and three sectors per cell, the voice channels are split into 21 groups,
so each sector has about $395/21\approx19$ channels.

Calls arrive at random, so a sector with 19 channels cannot carry 19\,erlangs of traffic without
blocking. The Erlang~B formula (Chapter~\ref{ch:multipleaccess}) gives the offered load at a
2\% blocking probability as $A\approx12.3$\,erlangs per sector, or about 37\,erlangs per cell.
A typical 1980s subscriber offered about 0.025\,erlang in the busy hour (a call of
about 90\,s each busy hour). One cell therefore served roughly $37/0.025\approx1500$
subscribers. A metropolitan system with a hundred cells served about 150\,000, a big improvement
on IMTS's two thousand, but only because many cells had been built. Halving the cell radius
quadrupled the subscribers and the number of cell sites. Note also the trunking inefficiency of
splitting the channels into 21 small groups: one 395-channel pool would carry about 380\,erlangs
at 2\% blocking, whereas 21 groups of 19 carry only about 260. That waste is a direct cost of
reuse-7 and sectorisation, and is one of the reasons 2G designers wanted tighter reuse.
\end{worked}

\subsection{NMT, TACS and the rest of the analog zoo}

AMPS was not first. NTT opened a cellular network in Tokyo in December 1979. The
\term{Nordic Mobile Telephone} system (NMT), designed jointly by the telecommunications
administrations of Denmark, Finland, Norway and Sweden, opened in 1981 at 450\,MHz with
25\,kHz channels, later joined by NMT-900. NMT's lasting contribution was organisational
rather than technical: it was the first system to offer \term{roaming} across national
borders, because its designers had insisted on a common specification for four countries. Its
open specification also spawned a Nordic handset industry; Nokia and Ericsson both learnt the
trade on NMT. The UK adopted a modified AMPS, \term{TACS} (Total Access Communication System),
with 25\,kHz channels at 900\,MHz, launched by Vodafone and Cellnet in 1985. West Germany had
C-Netz (450\,MHz), France had Radiocom~2000, Italy had RTMS, and many other countries had
their own systems. A traveller in 1988 Europe needed a different phone in almost every
country. That fragmentation was the political motivation for GSM.

\subsection{Why 1G had to go: capacity, privacy and cloning}

By the late 1980s analog systems were running out of capacity in large cities, and cell
splitting was getting expensive: below about 1\,km radius, sites were hard to find, handoffs
became frequent and co-channel interference from tall sites was hard to control. Digital
speech coding offered a way out: a digital channel could be shared by several users, tolerate
lower signal-to-interference ratios with error-control coding, and be encrypted.

The security failures were worse. An analog FM conversation could be heard by anyone with a
scanner tuned to the 870--890\,MHz band, and in the US such scanners were widely sold until
Congress banned the manufacture of scanners able to receive the cellular bands in the early
1990s. More damaging was \term{cloning}. An AMPS phone identified itself with its
\term{mobile identification number} (MIN, the phone number) and a 32-bit
\term{electronic serial number} (ESN) burnt in at the factory. Both were sent in the clear on
the reverse control channel every time the phone registered or made a call. A criminal with a
receiver near a busy motorway could harvest thousands of MIN/ESN pairs an hour, reprogram
other phones with them, and sell calls billed to the victims. The US industry's losses to
cloning fraud were estimated in the hundreds of millions of dollars a year by the mid-1990s.
Operators responded with RF fingerprinting (identifying individual transmitters by their
turn-on transients), personal identification numbers and, finally, a cryptographic
challenge--response protocol (the CAVE algorithm with a secret A-key) added to AMPS and its
digital successors. GSM's designers had the same problem in view, and their answer---a
secret key in a smart card that never leaves it---is still the basis of mobile security.

\begin{pitfall}[Identity is not authentication]
AMPS confused \emph{identification} (``I am MIN 312-555-0100, ESN 0x8A3F\ldots'') with
\emph{authentication} (``and here is proof''). Any identifier that is transmitted can be copied.
Authentication requires a secret that is never transmitted, used to answer a fresh challenge.
The same mistake recurs in modern systems wherever a static identifier is trusted: MAC-address
filtering in Wi-Fi, the IMSI sent in the clear in GSM and LTE attach procedures (which enables
the ``IMSI catcher''), or a device serial number used as a password in an IoT product. The
fix was the same each time: a challenge--response protocol based on a secret key, and, in 5G,
encryption of the permanent identity itself (Section~\ref{sec:ch21:5gc}).
\end{pitfall}

%======================================================================
\section{Second generation: GSM in depth}
\label{sec:ch21:gsm}
%======================================================================

The \term{Global System for Mobile communications} (GSM) is the most successful communication
standard ever written. At its peak, around 2012, roughly four billion people used it, and in
2025 it still carries voice, SMS and machine-to-machine traffic in much of the world. It is
also the best system to study first, because almost every idea in later generations appears
in it in a simpler form: a TDMA frame with a known training sequence, an adaptive equaliser,
slow frequency hopping, discontinuous transmission, power control, mobile-assisted handover,
a subscriber identity module, challenge--response authentication and a layered,
standardised core network with open interfaces between vendors.

\begin{history}[Groupe Spécial Mobile and the 1987 MoU]
In 1982 the Conference of European Postal and Telecommunications Administrations (CEPT) set up
a study group, the \emph{Groupe Spécial Mobile}, to specify a pan-European mobile system in the
900\,MHz band that the European Commission was preparing to reserve. Under its chairman Thomas
Haug of the Swedish administration, the group compared eight candidate radio systems in field
trials in Paris in 1986: several narrowband and wideband TDMA designs, frequency-hopping CDMA
and a narrowband FDMA system. In February 1987, at a meeting in Madeira, the group chose
narrowband TDMA with GMSK modulation, a 200\,kHz carrier and eight time slots, a compromise
between the French--German preference for a wideband system and the Nordic preference for a
narrowband one. On 7~September 1987, operators from thirteen countries signed a Memorandum of
Understanding in Copenhagen committing to open GSM networks by 1991. The MoU, not the technical
specification, is often called GSM's real innovation: it created a guaranteed market, so
manufacturers could invest in chips and handsets before a single network existed. In 1989
responsibility for the specifications moved to the newly formed European Telecommunications
Standards Institute (ETSI). The first GSM call was made in Finland on 1~July 1991 on the
Radiolinja network, and the acronym was soon redefined as Global System for Mobile
communications as networks opened on every continent.
\end{history}

\subsection{Design goals}

GSM's designers set themselves a list that still reads well: international roaming with a
single subscription; better speech quality and spectral efficiency than analog; security
against eavesdropping and fraud; low-cost handsets and infrastructure through open
interfaces, so that operators could buy base stations from one vendor and switches from
another; compatibility with the Integrated Services Digital Network (ISDN), then expected to
be the future of fixed telephony; and support for data and new services. Every major choice
below follows from these goals together with the radio environment at 900\,MHz: delay
spreads up to about 10--20\,$\mu$s in hilly terrain, vehicle speeds up to about 250\,km/h, and
cells up to 35\,km in radius.

\subsection{The radio interface: carriers, slots and GMSK}

The original GSM band (now called P-GSM~900) is 890--915\,MHz for the mobile transmit (uplink)
and 935--960\,MHz for the base transmit (downlink), with 45\,MHz duplex spacing and
124 carriers of 200\,kHz in each direction. Later allocations added 880--890/925--935\,MHz
(E-GSM), DCS~1800 (1710--1785/1805--1880\,MHz, 374 carriers) and PCS~1900 in the Americas
(3GPP TS 45.005). Each carrier is shared by eight users in time: GSM is a hybrid
\term{FDMA/TDMA} system.

The bit rate on a carrier is
\begin{equation}
R=\frac{13\,\mathrm{MHz}}{48}=270.833\ \mathrm{kb/s},\qquad T_b=\frac{48}{13}\,\mu\mathrm{s}\approx3.69\,\mu\mathrm{s},
\label{eq:ch21:gsmrate}
\end{equation}
derived from the 13\,MHz master clock found in every GSM device. The modulation is
\term{Gaussian minimum-shift keying} (GMSK) with a bandwidth--time product $BT=0.3$
(Chapter~\ref{ch:modulation}): continuous-phase binary FSK with modulation index $1/2$, in
which each bit's frequency pulse has been smoothed by a Gaussian filter. The smoothing spreads
each pulse over about three bit periods (Figure~\ref{fig:ch21_gmsk}b), which costs a little
intersymbol interference but makes the spectrum fall much faster than plain MSK. GMSK was
chosen because it has a \emph{constant envelope}: the handset's power amplifier can run in
saturation, at high efficiency, without spectral regrowth (Chapter~\ref{ch:sdr}). In 1987 a
linear amplifier for a handset was expensive and power-hungry; a class-C amplifier was cheap.
The price is spectral efficiency: one bit per symbol, 1.35\,b/s/Hz gross.

\mdcfig{ch21_gmsk}{(a) Power spectra of MSK and GMSK at the GSM bit rate, with an approximate
outline of the GSM modulation spectrum mask (after TS 45.005; the real mask depends on power
class and measurement bandwidth). With $BT=0.3$ the spectrum is about 30\,dB down at the
adjacent-channel centre (200\,kHz) and below $-60$\,dB at 400\,kHz; MSK's sidelobes would
violate the mask. (b) The frequency pulse of one bit: MSK's is rectangular, GMSK's is a
Gaussian-smoothed pulse spread over about three bit periods, which is the price paid in ISI
for the compact spectrum.}

\subsection{The time structure}

Figure~\ref{fig:ch21_gsm_frames} shows GSM's time hierarchy (3GPP TS 45.002). A
\term{timeslot} lasts 156.25 bit periods, 576.9\,$\mu$s. Eight timeslots make a
\term{TDMA frame} of 4.615\,ms ($=60/13$\,ms). A user with a full-rate traffic channel owns one
timeslot in every frame. Frames are grouped into \term{multiframes}: 26 frames (120\,ms) for
traffic channels, 51 frames (235.4\,ms) for the broadcast and common control channels. The two
lengths were chosen to be coprime so that a mobile on a traffic channel, which has an idle
frame in every 26, slides past every position of the 51-multiframe of a neighbouring cell and
can find its synchronisation burst within a known time. A superframe of
$26\times51$ frames (6.12\,s) contains every combination, and a hyperframe of 2048 superframes
(3\,h\,28\,min\,53.76\,s) sets the period of the frame number used by the cipher and the
frequency-hopping sequence.

\mdcfig[0.95]{ch21_gsm_frames}{The GSM time hierarchy, from the 3.5-hour hyperframe down to
the normal burst. A full-rate traffic user owns one timeslot per TDMA frame; in each
26-multiframe, 24 frames carry traffic, one carries the slow associated control channel
(SACCH, ``S'') and one is idle (``I''), which the mobile uses to measure neighbouring cells.
The normal burst places its 26-bit training sequence in the middle, so that the channel
estimate is valid for both halves of the data.}

\subsubsection{Bursts}

What is sent in a timeslot is a \term{burst}. The \term{normal burst} carries two 57-bit data
fields, a 26-bit \term{training sequence} in the middle (a \emph{midamble}), two
\term{stealing flags} that say whether the data fields have been borrowed for urgent
signalling, three tail bits at each end and a guard period of 8.25 bit periods (about
30\,$\mu$s) during which the transmitter ramps its power up and down. The design choices are
worth studying.
\begin{itemize}
\item \textbf{The midamble.} The channel changes during a burst at high speed: at 250\,km/h and
900\,MHz the maximum Doppler is about 210\,Hz, a coherence time of a couple of milliseconds,
not much more than the 0.58\,ms burst. Placing the training sequence in the middle keeps every
data bit within about 0.3\,ms of the channel estimate.
\item \textbf{The training sequences.} There are eight \term{training sequence codes} (TSC),
each a 16-bit core extended cyclically by five bits on each side. Within $\pm5$ bits of the
peak their autocorrelation is ideal (zero off-peak), so correlating the received midamble with
the known sequence gives the channel impulse response directly over a window of about five
bits, roughly 18\,$\mu$s. Neighbouring co-channel cells are given different TSCs, so the
estimate favours the wanted signal.
\item \textbf{The equaliser.} With the channel estimated, a Viterbi (MLSE) equaliser with a
four- or five-bit memory (16--32 states) detects the data (Chapter~\ref{ch:equalization}).
The GSM specification requires good performance on its ``hilly terrain'' channel model with
echoes delayed by about 15--20\,$\mu$s, so the equaliser is mandatory, and it was the most
demanding block in a 1991 handset. GMSK's own ISI from the Gaussian filter is equalised in
the same trellis.
\item \textbf{The guard period.} 30\,$\mu$s covers the power ramps but not the round-trip
propagation delay to a distant mobile, which is up to about 230\,$\mu$s at 35\,km. That is
handled by timing advance, below.
\end{itemize}
Four other burst types complete the set. The \term{frequency correction burst} is 142 zeros
which, in GMSK, produce a pure tone 67.7\,kHz ($=R/4$) above the carrier, easy to find with a
narrow filter during cell search. The \term{synchronisation burst} carries a long 64-bit
training sequence and the base station identity code and reduced frame number. The
\term{access burst}, used for the very first transmission from a mobile that does not yet know
its distance, has a guard period of 68.25 bit periods (252\,$\mu$s), enough for a round trip
of about 35\,km, the maximum cell radius of standard GSM. The \term{dummy burst} fills unused
slots on the broadcast carrier, which must always be transmitted at constant power so that
mobiles in neighbouring cells can measure it.

\subsubsection{Timing advance and half-duplex operation}

A mobile at distance $d$ sees the base's frames late by $d/c$ and its own bursts arrive at the
base a further $d/c$ late. To keep bursts from different mobiles from overlapping, the base
measures the delay of each mobile's bursts and commands a \term{timing advance} (TA) of
0 to 63 bit periods; each step is 3.69\,$\mu$s of round trip, about 553\,m of distance, so
63 steps cover about 35\,km. The uplink frame is also delayed by three timeslots relative to
the downlink, so that a phone using timeslot 2 receives in slot 2, transmits three slots
later and never transmits and receives at the same instant. The phone needs no duplex
filter to separate simultaneous transmission and reception; a transmit/receive switch is
enough. This was a significant saving in 1991 and remains a reason why GSM chips are
cheap.

\subsection{Logical channels}

On top of the physical structure GSM defines \term{logical channels}, listed in
Table~\ref{tab:ch21:gsmlogical}. Traffic channels carry speech or data. Control channels come
in three families: \emph{broadcast} channels, which every mobile in a cell listens to;
\emph{common control} channels, shared by all mobiles for paging and access; and
\emph{dedicated} control channels, assigned to one mobile. The broadcast and common channels
live on timeslot~0 of one special carrier per cell, the \term{BCCH carrier}, which transmits
continuously at constant power and does not hop; that is the carrier mobiles measure for
handover and cell reselection.

\begin{table}[htbp]\centering\small
\caption{GSM logical channels (3GPP TS 45.002).}
\label{tab:ch21:gsmlogical}
\begin{tabularx}{\linewidth}{@{}l l X@{}}
\toprule
Channel & Direction & Purpose\\
\midrule
TCH/F, TCH/H & both & Full-rate (22.8\,kb/s gross) and half-rate (11.4\,kb/s) traffic: speech or circuit data\\
FCCH & down & Frequency correction bursts: a tone for coarse frequency and slot timing\\
SCH & down & Synchronisation: base station identity code (BSIC), TDMA frame number\\
BCCH & down & System information: cell identity, neighbour list, access parameters\\
PCH & down & Paging: ``mobile X, you have a call''\\
RACH & up & Random access: slotted-ALOHA access bursts requesting a channel\\
AGCH & down & Access grant: assigns a dedicated channel after a RACH request\\
SDCCH & both & Stand-alone dedicated control: call set-up, location updates, SMS when idle\\
SACCH & both & Slow associated control: power control, timing advance, measurement reports (every 480\,ms)\\
FACCH & both & Fast associated control: handover commands, by stealing traffic bursts\\
\bottomrule
\end{tabularx}
\end{table}

A mobile-originated call shows the channels in action. The mobile sends an access burst on
the RACH with a random reference; the network replies on the AGCH with an SDCCH assignment and
an initial timing advance; on the SDCCH the mobile and network authenticate, start ciphering
and exchange call set-up messages; the network then assigns a TCH, and the conversation runs
with the SACCH alongside it, carrying the mobile's measurements of its serving cell and up to
six neighbours every 480\,ms. If a handover is needed, the command arrives on the FACCH, which
steals a 20\,ms speech frame; the stealing flags tell the receiver to decode the burst as
signalling, and the listener hears a short gap concealed by the speech decoder.

\subsection{Speech coding and channel coding}

The original \term{full-rate} speech codec, regular-pulse-excitation long-term-prediction
(RPE-LTP, GSM~06.10), produces 260 bits every 20\,ms, 13\,kb/s. The bits are not equally
important, and GSM was one of the first systems to protect them unequally
(Chapter~\ref{ch:sourcecoding} explains why):
\begin{itemize}
\item 50 \emph{class 1a} bits (filter coefficients, block amplitudes, pitch parameters), whose
errors are audible as clicks, get a 3-bit CRC; if the CRC fails, the whole frame is discarded
and the decoder repeats and attenuates the previous one (\emph{bad-frame handling});
\item 132 \emph{class 1b} bits are convolutionally coded but not checked;
\item 78 \emph{class 2} bits are sent unprotected.
\end{itemize}
The class 1 bits, the 3 parity bits and 4 tail bits ($182+3+4=189$) pass through a rate-1/2,
constraint-length-5 convolutional code (generators $1+D^3+D^4$ and $1+D+D^3+D^4$,
Chapter~\ref{ch:classiccodes}), giving 378 bits, to which the 78 class 2 bits are added: 456
coded bits per 20\,ms, 22.8\,kb/s. The 456 bits are interleaved diagonally over eight
consecutive bursts, $8\times57=456$, so each burst carries halves of two speech frames. A
fade that wipes out one burst destroys only one eighth of each codeword, which the Viterbi
decoder can usually repair.

\begin{worked}[GSM's rates, from the slot to the speech codec]
\emph{Gross rate per slot.} The carrier runs at 270.833\,kb/s, shared by 8 slots:
$270.833/8=33.85$\,kb/s per slot, including tails, training, flags and guard.

\emph{Payload per slot.} Each normal burst carries $2\times57=114$ data bits every TDMA frame of
4.615\,ms: $114/4.615\,\mathrm{ms}=24.7$\,kb/s. The burst overhead is therefore
$1-114/156.25=27\%$.

\emph{Traffic channel.} In each 120\,ms 26-multiframe, 24 frames carry traffic: $24\times114$ bits
per 120\,ms $=22.8$\,kb/s, exactly the coded speech rate computed above. The 26th frames carry
the SACCH ($114$ bits per 120\,ms, 950\,b/s gross) and the idle frame.

\emph{Net speech rate.} 13\,kb/s of RPE-LTP. The channel-coding rate is $13/22.8=0.57$.

\emph{Spectral efficiency.} Eight 13\,kb/s voices in 200\,kHz is 0.52\,b/s/Hz per carrier
before reuse. With a reuse of 4 sites and 3 sectors (12 carrier groups), each sector gets a
twelfth of the band: about 0.043\,b/s/Hz per sector, or one simultaneous call for every
$12\times25=300$\,kHz of spectrum. AMPS, with 30\,kHz channels in 21 groups, needed
$21\times30=630$\,kHz per call per sector. Counted that way, GSM with 4/12 reuse is about twice
as efficient as AMPS, and with tighter reuse, frequency hopping and half-rate or AMR codecs
roughly 3--6 times. That is
where 2G's capacity came from: not from the modulation, which is no denser, but from tolerance
of interference (a C/I requirement of about 9\,dB instead of 18\,dB) that allowed tighter reuse.
\end{worked}

\subsubsection{From RPE-LTP to AMR}

Speech coding improved through GSM's life. The half-rate codec (VSELP, 5.6\,kb/s, 1995)
doubled capacity at some loss of quality. The enhanced full-rate codec (EFR, ACELP at
12.2\,kb/s, 1996) gave close to wireline quality in the full-rate channel. The
\term{Adaptive Multi-Rate} (AMR) codec (3GPP, 1999) provides eight modes from 4.75 to
12.2\,kb/s and lets the network choose, every 40\,ms or so, how to split the gross channel
between speech and protection: a user near the cell edge gets a low-rate codec with strong
coding, a user near the base gets 12.2\,kb/s with light coding. AMR can also run in a
half-rate channel. AMR-WB (G.722.2, 6.6--23.85\,kb/s) brought wideband ``HD voice'' (50--7000\,Hz)
to GSM and UMTS. AMR is the first appearance in this chapter of \emph{link adaptation}, the
principle that later dominates HSPA, LTE and NR.

\subsection{Frequency hopping, DTX and power control}

GSM supports \term{slow frequency hopping}: a traffic channel changes carrier every TDMA
frame, 217 times a second, following either a cyclic or a pseudo-random sequence over a list
of carriers allocated to the cell (Chapter~\ref{ch:spreadspectrum}). Each 20\,ms speech frame is
interleaved over eight bursts on eight different frequencies, which gives two kinds of gain.
\emph{Frequency diversity} helps a slow-moving or stationary user, whose channel would otherwise
sit in a deep fade for seconds: with hopping, a fade on one carrier destroys only one burst of
eight. \emph{Interference diversity} helps everyone: with uncorrelated hopping sequences in
neighbouring cells, the interference on each burst comes from a different random user, so the
codeword sees average rather than worst-case interference. Hopping is the main reason GSM
networks could move from 4/12 reuse to 3/9 and even 1/3 or 1/1 with fractional loading in the
late 1990s. It was an elegant way of getting part of CDMA's interference averaging inside a
TDMA system.

\textbf{Discontinuous transmission}\index{discontinuous transmission} (DTX) exploits the fact that a person in a conversation
talks less than half the time. A voice activity detector stops transmission during silences,
except for occasional silence-descriptor frames from which the receiver synthesises
\emph{comfort noise} so that the line does not sound dead. DTX roughly halves the average
interference and saves battery. \textbf{Power control}\index{power control}, in 2\,dB steps every SACCH period
(480\,ms) over a range of about 30\,dB, keeps each mobile's power no higher than needed, which
again reduces interference and saves battery. A GSM~900 handset (power class~4) transmits at
most 2\,W during its burst, an average of 250\,mW because it transmits in only one slot in
eight.

\begin{inpractice}[The 217\,Hz buzz]
Anyone who has left a GSM phone next to a loudspeaker or a car radio has heard a buzzing
``dit-dit-dit-dit'': that is the 217\,Hz TDMA frame rate, rectified by audio amplifiers that
were never designed to see a 2\,W, 900\,MHz burst switching on and off 217 times a second. The
same pulsed envelope caused interference in hearing aids and early medical devices and led to
the hearing-aid compatibility ratings still printed on phones. CDMA, WCDMA and OFDMA phones
transmit continuously during a call (or at least with much less regular gating), and the buzz
disappeared with 2G. It is a reminder that the envelope of a signal matters to equipment that
does not even demodulate it.
\end{inpractice}

\subsection{Network architecture}

GSM's architecture (Figure~\ref{fig:ch21_gsm_arch}) divides the network into a
\term{base station subsystem} (BSS) and a \term{network and switching subsystem} (NSS), with
standardised interfaces between them.
\begin{itemize}
\item The \term{mobile station} (MS) is the \emph{mobile equipment} plus the
\term{subscriber identity module} (SIM), a smart card holding the subscriber's identity (IMSI)
and secret key.
\item The \term{base transceiver station} (BTS) contains the radios and handles the physical
layer: modulation, equalisation, channel coding, ciphering and frequency hopping.
\item The \term{base station controller} (BSC) manages tens to hundreds of BTSs over the
\emph{Abis} interface: it allocates channels, runs power control and decides on handovers
between its own cells. A transcoder unit (TRAU) converts between the 13\,kb/s codec and the
64\,kb/s PCM of the fixed network.
\item The \term{mobile switching centre} (MSC), a telephone exchange extended for mobility, connects
calls over the \emph{A} interface and to the public telephone network through a gateway MSC.
The \term{visitor location register} (VLR) holds data on the subscribers currently in its area.
\item The \term{home location register} (HLR) is the master database of the subscribers of a
network: their identities, subscribed services and the VLR that currently serves them. The
\term{authentication centre} (AuC) holds each subscriber's secret key and generates
authentication vectors. The equipment identity register (EIR) lists stolen handsets.
\end{itemize}
The core elements talk to each other with the Mobile Application Part (MAP) of Signalling
System No.~7, the same signalling network that set up fixed telephone calls
(Chapter~\ref{ch:wireline}). Roaming is a MAP conversation between the visited network's VLR
and the home network's HLR: ``a subscriber of yours has turned up here; is she allowed to
use my network, and please send me her authentication vectors.''

\begin{figure}[htbp]\centering
\begin{tikzpicture}[font=\sffamily\scriptsize,
  box/.style={draw=navy,fill=navy!6,thick,rounded corners=1pt,minimum height=0.8cm,minimum width=1.3cm,align=center},
  db/.style={draw=accent,fill=accent!6,thick,rounded corners=1pt,minimum height=0.7cm,minimum width=1.1cm,align=center},
  pk/.style={draw=practc,fill=practc!8,thick,rounded corners=1pt,minimum height=0.8cm,minimum width=1.3cm,align=center},
  ln/.style={thick,navy}, if/.style={font=\sffamily\tiny,text=accent,midway,above}]
\node[box] (ms) at (0,0) {MS\\(ME + SIM)};
\node[box] (bts1) at (2.4,0.9) {BTS};
\node[box] (bts2) at (2.4,-0.9) {BTS};
\node[box] (bsc) at (4.8,0) {BSC\\+ TRAU};
\node[box] (msc) at (7.4,0.9) {MSC/VLR};
\node[box] (gmsc) at (10.1,0.9) {GMSC};
\node[draw=gray,thick,ellipse,minimum width=1.4cm,minimum height=0.8cm] (pstn) at (12.4,0.9) {PSTN};
\node[db] (hlr) at (8.75,2.4) {HLR/AuC};
\node[db] (eir) at (6.6,2.4) {EIR};
\node[db] (smsc) at (10.9,2.4) {SMSC};
\node[pk] (sgsn) at (7.4,-0.95) {SGSN};
\node[pk] (ggsn) at (10.1,-0.95) {GGSN};
\node[draw=gray,thick,ellipse,minimum width=1.4cm,minimum height=0.8cm] (inet) at (12.4,-0.95) {Internet};
\draw[ln,dashed] (ms) -- node[if]{Um} (bts1); \draw[ln,dashed] (ms) -- (bts2);
\draw[ln] (bts1) -- node[if]{Abis} (bsc); \draw[ln] (bts2) -- (bsc);
\draw[ln] (bsc) -- node[if]{A} (msc); \draw[ln,practc] (bsc) -- node[if,below]{Gb} (sgsn);
\draw[ln] (msc) -- (gmsc); \draw[ln] (gmsc) -- (pstn);
\draw[ln,practc] (sgsn) -- node[if]{Gn} (ggsn); \draw[ln,practc] (ggsn) -- node[if]{Gi} (inet);
\draw[ln,accent!70,dotted] (msc) -- node[if,left]{MAP} (hlr); \draw[ln,accent!70,dotted] (gmsc) -- (hlr);
\draw[ln,accent!70,dotted] (msc) -- (eir); \draw[ln,accent!70,dotted] (sgsn.north east) -- (hlr);
\draw[ln,accent!70,dotted] (gmsc) -- (smsc);
\node[draw=navy!40,dashed,rounded corners,fit=(bts1)(bts2)(bsc),inner sep=4pt,label={[font=\sffamily\tiny,text=navy]below:base station subsystem}] {};
\node[draw=navy!40,dashed,rounded corners,fit=(msc)(gmsc)(hlr)(eir)(smsc),inner sep=4pt,label={[font=\sffamily\tiny,text=navy]above:network and switching subsystem (circuit)}] {};
\node[draw=practc!60,dashed,rounded corners,fit=(sgsn)(ggsn),inner sep=4pt,label={[font=\sffamily\tiny,text=practc]below:GPRS packet core (2.5G)}] {};
\end{tikzpicture}
\caption{GSM network architecture with the GPRS packet core added in the late 1990s. Solid
lines carry traffic and signalling; dotted lines are SS7/MAP signalling to the databases.
Standardised interfaces (Um, Abis, A, Gb) let operators buy each element from a different
vendor, one of GSM's founding goals.}
\label{fig:ch21_gsm_arch}
\end{figure}

\subsection{Security: the SIM, authentication and the A5 ciphers}

GSM's security design (Figure~\ref{fig:ch21_gsm_auth}) rests on a 128-bit secret key $K_i$
shared between the SIM and the AuC and never transmitted. When the network wants to
authenticate a mobile, the AuC generates a random 128-bit challenge RAND and computes two
functions of $K_i$ and RAND: a 32-bit signed response SRES (algorithm A3) and a 64-bit
cipher key $K_c$ (algorithm A8). The triplet (RAND, SRES, $K_c$) is sent to the serving
MSC/VLR, which sends RAND to the mobile. The SIM computes SRES and $K_c$ itself and returns
SRES; if it matches, the subscriber is authentic. Both sides now hold $K_c$ without its ever
having crossed the air, and the BTS and the handset encrypt each burst with the stream cipher
\term{A5/1} (or the export variant A5/2, or later the KASUMI-based A5/3), keyed by $K_c$ and
the TDMA frame number. A5/1 is built from three short linear-feedback shift registers of
lengths 19, 22 and 23 bits with irregular, majority-rule clocking, and it produces 114
keystream bits for each direction in each frame. To protect location privacy the network
assigns a \term{temporary mobile subscriber identity} (TMSI) after the first registration, so
the permanent IMSI need not be sent again.

\begin{figure}[htbp]\centering
\begin{tikzpicture}[font=\sffamily\scriptsize,
  box/.style={draw=navy,fill=navy!6,thick,rounded corners=1pt,minimum height=0.75cm,align=center},
  alg/.style={draw=accent,fill=accent!6,thick,minimum height=0.55cm,minimum width=0.9cm,align=center},
  arr/.style={-{Stealth[length=2mm]},thick,navy}]
\node[box,minimum width=3.6cm,minimum height=3.3cm] (sim) at (0,0) {};
\node[font=\sffamily\bfseries\scriptsize,text=navy] at (0,1.4) {SIM (in the handset)};
\node[box,minimum width=4.4cm,minimum height=3.3cm] (auc) at (8.2,0) {};
\node[font=\sffamily\bfseries\scriptsize,text=navy] at (8.2,1.4) {AuC / HLR (home network)};
\node[alg] (a3s) at (-0.8,0.2) {A3};
\node[alg] (a8s) at (0.8,0.2) {A8};
\node (kis) at (0,0.95) {$K_i$ (128 bits, never leaves)};
\node[alg] (a3n) at (7.4,0.2) {A3};
\node[alg] (a8n) at (9.0,0.2) {A8};
\node (kin) at (8.2,0.95) {$K_i$ \quad RAND generator};
\node (sress) at (-0.8,-0.9) {SRES};
\node (kcs) at (0.8,-0.9) {$K_c$ (64 bits)};
\node (sresn) at (7.4,-0.9) {SRES};
\node (kcn) at (9.0,-0.9) {$K_c$};
\draw[arr] (kis) -- (a3s); \draw[arr] (kis) -- (a8s); \draw[arr] (kin.south west) -- (a3n); \draw[arr] (kin) -- (a8n);
\draw[arr] (a3s) -- (sress); \draw[arr] (a8s) -- (kcs);
\draw[arr] (a3n) -- (sresn); \draw[arr] (a8n) -- (kcn);
\node[box,minimum width=2.4cm,fill=practc!8,draw=practc] (vlr) at (4.1,-0.9) {MSC/VLR\\compare SRES};
\draw[arr,accent] (6.0,0.6) -- node[above,font=\sffamily\tiny]{1. triplets (RAND, SRES, $K_c$)} (4.1,0.6) -- (vlr.north);
\draw[arr] (vlr.north west) ++(0,0.0) -- node[above,sloped,font=\sffamily\tiny]{2. RAND} (1.8,0.2);
\draw[arr] (1.8,-0.9) -- node[below,font=\sffamily\tiny]{3. SRES} (vlr.west);
\node[draw=gray,dashed,rounded corners,align=center,font=\sffamily\tiny] at (4.1,-2.15)
  {4. BTS and handset encrypt each burst with A5/1 (or A5/3), key $K_c$, frame number as IV};
\end{tikzpicture}
\caption{GSM authentication and key agreement. The secret $K_i$ stays in the SIM and the AuC;
only the random challenge and the 32-bit response cross the air, and both ends derive the
cipher key $K_c$ independently. The network authenticates the mobile, but the mobile never
authenticates the network.}
\label{fig:ch21_gsm_auth}
\end{figure}

\begin{pitfall}[How GSM's security was broken]
GSM's security was a large advance on AMPS, but nearly every weakness that cryptographers would
later exploit was a design decision made for cost or politics, and each is now a textbook
lesson.
\begin{itemize}
\item \emph{Secret algorithms.} A3/A8 and A5 were confidential. Once reverse-engineered, they
were broken quickly. In 1998 Briceno, Goldberg and Wagner showed that COMP128-1, the A3/A8
example algorithm many operators used, leaked $K_i$ after roughly $10^5$ chosen challenges,
allowing SIM cloning; COMP128-1 also set ten bits of $K_c$ to zero, weakening A5 to an
effective 54-bit key. A5/2 was broken almost immediately after publication in 1999.
\item \emph{Coding before ciphering.} GSM encrypts \emph{after} error-control coding, so the
ciphertext obeys known linear parity relations. Barkan, Biham and Keller (2003) used this to
break A5/2 in real time from ciphertext alone, and showed active attacks that tricked a phone
into using A5/2 and so recovered the $K_c$ also used with A5/1.
\item \emph{A5/1 itself.} Time--memory trade-off attacks (Biryukov, Shamir and Wagner, 2000) and,
from 2009--10, publicly computed rainbow tables of about 2\,TB made A5/1 decryptable on a PC
given a few known-plaintext bursts, which GSM's predictable signalling messages supply.
\item \emph{One-way authentication.} The network authenticates the phone, not vice versa. A
fake base station (an \term{IMSI catcher}) can ask the phone for its IMSI, tell it not to
encrypt, and relay its calls.
\item \emph{No integrity protection} of signalling and encryption only on the radio link.
\end{itemize}
UMTS fixed most of these with mutual authentication (AKA), integrity protection, published
algorithms (KASUMI, later SNOW~3G and AES) and 128-bit keys; LTE and 5G kept that model.
Kerckhoffs's principle---assume the attacker knows the algorithm---was learnt the hard way.
\end{pitfall}

\subsection{SMS: the accidental success}

The \term{Short Message Service} was an afterthought. The idea of sending short text messages
over the signalling channels of the network was developed in the mid-1980s by Friedhelm
Hillebrand and Bernard Ghillebaert in the GSM standardisation work, and the limit of 160
characters came from the 140 octets available in a signalling message at 7 bits per
character. Hillebrand later recounted that he had checked, by typing out typical postcards and
telex messages, that 160 characters was enough for most of what people wanted to say. The
first SMS, ``Merry Christmas'', was sent from a PC to a phone on the Vodafone UK network on
3~December 1992. Because messages travel on the SDCCH or SACCH and are stored and forwarded
by a short-message service centre, SMS costs the network almost nothing: it uses capacity
that already exists for signalling. Operators priced it as if it were valuable, teenagers
discovered it once prepaid tariffs and inter-network messaging arrived in the late 1990s, and
for a decade SMS was the most profitable service in telecommunications, with several trillion
messages a year by around 2010. The lesson, repeated by the mobile internet and by apps,
is that the most important use of a network is rarely the one its designers planned.

\begin{inpractice}[GSM's long goodbye]
Operators are switching off 2G and 3G networks to reuse the spectrum for 4G and 5G. AT\&T in
the US switched off its GSM network at the start of 2017 and its 3G network in 2022; many
European operators have chosen to keep GSM longer than 3G, because GSM serves roaming
visitors, old phones, alarm panels, smart meters and cars with emergency-call modules, and
because a single GSM carrier fits in a 200\,kHz sliver between LTE carriers. GSM-R, the railway
version, carries train signalling (ETCS level 2) across Europe and is being replaced by the
5G-based FRMCS from around 2030. Engineers planning an IoT product with a 15-year life should
check each target country's switch-off plans before choosing a 2G or 3G module.
\end{inpractice}

%======================================================================
\section{The other second-generation systems}
\label{sec:ch21:other2g}
%======================================================================

GSM did not have the world to itself. In the US the FCC did not mandate a digital standard,
operators wanted to upgrade in their existing AMPS spectrum, and two incompatible digital
systems emerged; Japan built its own.

\subsection{D-AMPS (IS-54/IS-136) and PDC}

The US industry association (CTIA) chose TDMA as its digital standard in 1988, and the result,
IS-54 (1990), later IS-136, was called \term{D-AMPS} or simply ``TDMA''. Its constraint was
backward compatibility: it had to fit in a 30\,kHz AMPS channel so that operators could convert
channels one at a time and dual-mode phones could fall back to AMPS. Each 30\,kHz carrier
carries 48.6\,kb/s of $\pi/4$-DQPSK at 24.3\,kBd with root-raised-cosine filtering
($\beta=0.35$) in 40\,ms frames of six slots; a full-rate user takes two slots, so three users
share a channel that carried one under AMPS. The speech codec was 7.95\,kb/s VSELP. A
three-fold capacity gain looked good, but IS-136 never had GSM's richness of features or its
economies of scale; AT\&T Wireless and Cingular moved their TDMA networks to GSM in the early
2000s. $\pi/4$-DQPSK, a linear modulation, required a linear power amplifier but gave twice
GMSK's bits per hertz; the same trade-off reappears in EDGE.

Japan's \term{Personal Digital Cellular} (PDC, 1993) was similar: 25\,kHz carriers, three-slot
TDMA, $\pi/4$-DQPSK at 42\,kb/s. It was used only in Japan, but it carried NTT DoCoMo's
\emph{i-mode} service (1999), the first mass-market mobile internet, with tens of millions of
users sending email and browsing compact web pages years before the smartphone.

\subsection{IS-95 cdmaOne}

The radical alternative was Qualcomm's \term{code-division multiple access} system,
standardised as IS-95 in 1993 and marketed as \term{cdmaOne}. Its physical layer---direct-sequence
spreading, Walsh codes, PN offsets, the rake receiver---is developed in
Chapter~\ref{ch:spreadspectrum}; here the system view matters.
\begin{itemize}
\item \textbf{Wideband carriers.} 1.25\,MHz carriers spread at 1.2288\,Mchip/s, so one CDMA
carrier replaces 41 AMPS channels. Every cell uses the same carrier: \term{universal frequency
reuse} ($N=1$). Cells are distinguished by offsets of a common short PN code (512 offsets of
64 chips), users by orthogonal 64-chip Walsh codes on the downlink and by a long PN code on
the uplink.
\item \textbf{Interference-limited capacity.} Capacity is set by total interference rather than
by the number of channels, so anything that reduces interference adds users: voice activity
(the variable-rate QCELP and EVRC vocoders transmit at an eighth of full rate in silences),
sectorisation, and above all tight power control.
\item \textbf{Power control.} Uplink closed-loop power control at 800\,Hz in 1\,dB steps solves
the near--far problem, in which a nearby phone would otherwise drown a distant one by 60\,dB
or more. Section~\ref{sec:ch21:3g} examines power control in WCDMA, which inherited the idea
at a faster rate.
\item \textbf{Soft handoff.} A mobile near a cell boundary communicates with two or three
cells at once, and the network combines the copies. There is no break in the call and the
diversity reduces the power needed.
\item \textbf{Synchronised base stations.} The short-PN-offset scheme requires all base stations
to share a time reference, provided by GPS receivers at every site; this became a political
and operational issue in countries wary of dependence on GPS.
\end{itemize}

\begin{history}[Qualcomm's CDMA bet]
Qualcomm was founded in San Diego in 1985 by Irwin Jacobs, Andrew Viterbi and five colleagues,
most of whom had worked together at Linkabit on spread-spectrum military and satellite
systems. In 1988--89 they proposed applying CDMA to cellular, at the moment the industry had
just chosen TDMA. Their claims of a capacity many times that of AMPS met loud scepticism;
critics argued that power control would never be accurate enough in a real fading channel and
that the system would collapse under its own interference. Qualcomm demonstrated a working
prototype in San Diego in late 1989, persuaded the industry to standardise it as IS-95 in
1993, and won early commitments from operators in Hong Kong, South Korea and the US. South
Korea's government chose CDMA as its national digital standard, and Korean manufacturers
became its first mass producers. Real CDMA capacity turned out lower than the most optimistic
early claims but clearly higher than D-AMPS. The deeper victory was intellectual: every 3G
standard used CDMA, and Qualcomm's patent portfolio made it the most important licensor in
mobile communications.
\end{history}

\begin{table}[htbp]\centering\small
\caption{The main 2G air interfaces.}
\label{tab:ch21:2g}
\footnotesize
\begin{tabularx}{\linewidth}{@{}l l l l l >{\raggedright\arraybackslash}X@{}}
\toprule
System & Carrier & Access & Modulation & Gross rate & Speech codecs (kb/s)\\
\midrule
GSM & 200\,kHz & TDMA, 8 slots & GMSK & 270.8\,kb/s & RPE-LTP 13; EFR 12.2; AMR 4.75--12.2\\
D-AMPS & 30\,kHz & TDMA, 3 users & $\pi/4$-DQPSK & 48.6\,kb/s & VSELP 7.95; ACELP 7.4\\
PDC & 25\,kHz & TDMA, 3 users & $\pi/4$-DQPSK & 42\,kb/s & VSELP 6.7\\
IS-95 & 1.25\,MHz & CDMA, $N=1$ & QPSK/OQPSK & 1.2288\,Mc/s & QCELP 8 and 13; EVRC\\
\bottomrule
\end{tabularx}
\end{table}

%======================================================================
\section{Generation 2.5: GPRS and EDGE}
\label{sec:ch21:gprs}
%======================================================================

GSM was designed for circuit-switched voice, and its data service was a circuit too: 9.6\,kb/s
(later 14.4) on a traffic channel, held and billed for the whole session, like a modem call.
High-speed circuit-switched data (HSCSD) bundled up to four slots. Neither suited bursty
internet traffic, in which a user downloads a page in a second and then reads it for a minute.

\subsection{GPRS: packets on a circuit network}

The \term{General Packet Radio Service} (GPRS), commercial from about 2000, added packet
switching to GSM without changing the physical layer. Timeslots can be configured as
\term{packet data channels} (PDCH), shared among many users by a MAC scheduler in the BSC
(the packet control unit). Data is sent in radio blocks of four bursts over four consecutive
frames (20\,ms). A phone may use several slots at once, up to the limit of its
\term{multislot class}; a common class-10 phone could receive on four slots and transmit on
two, five in total. Four coding schemes trade protection for rate: CS-1 to CS-4 give 9.05,
13.4, 15.6 and 21.4\,kb/s per slot. An RLC layer with selective-repeat ARQ makes the link
reliable, and on the uplink a 3-bit \emph{uplink state flag} in each downlink block says which
mobile may transmit in the next uplink block, a simple form of the dynamic uplink grant that
LTE and NR use. In the core, the \term{SGSN} (serving GPRS support node) handles mobility and
sessions, and the \term{GGSN} (gateway GPRS support node) is the IP router to the internet; a
user activates a \emph{PDP context} identified by an access point name (APN), the ancestor of
the LTE bearer and the NR PDU session. Packets are tunnelled between the two in the GPRS
tunnelling protocol, GTP, which survives in 5G as GTP-U on the N3 interface.

GPRS gave the first ``always-on'' mobile internet, but it was slow (tens of kilobits per
second in practice) and had round-trip times of 600\,ms to over a second, because a packet
waited for resource allocation in the BSC and then crossed several network nodes. A user felt
the latency more than the rate.

\subsection{EDGE: a smarter modulation and incremental redundancy}

\term{EDGE} (Enhanced Data rates for GSM Evolution, also EGPRS) kept GSM's 200\,kHz carrier,
270.833\,ksym/s symbol rate and burst structure but added \term{8-PSK}, three bits per symbol.
To avoid zero crossings of the envelope (which stress a power amplifier), the constellation is
rotated by $3\pi/8$ every symbol, the same trick as $\pi/4$-QPSK
(Chapter~\ref{ch:modulation}). Nine modulation and coding schemes give 8.8 to 59.2\,kb/s per
slot (Table~\ref{tab:ch21:edge}); with four slots, a phone reaches 236.8\,kb/s at the best.
\begin{table}[htbp]\centering\small
\caption{EGPRS modulation and coding schemes (3GPP TS 44.060), data rate per timeslot.
Schemes in the same family share a payload unit size, so a block can be resegmented and
retransmitted with a more robust scheme of the same family.}
\label{tab:ch21:edge}
\begin{tabular}{@{}l c c c | l c c c@{}}
\toprule
Scheme & Mod. & Family & kb/s & Scheme & Mod. & Family & kb/s\\
\midrule
MCS-1 & GMSK & C & 8.8 & MCS-5 & 8-PSK & B & 22.4\\
MCS-2 & GMSK & B & 11.2 & MCS-6 & 8-PSK & A & 29.6\\
MCS-3 & GMSK & A & 14.8 & MCS-7 & 8-PSK & B & 44.8\\
MCS-4 & GMSK & C & 17.6 & MCS-8 & 8-PSK & A & 54.4\\
 & & & & MCS-9 & 8-PSK & A & 59.2\\
\bottomrule
\end{tabular}
\end{table}

EDGE's real intellectual content is in its link layer. \textbf{Link adaptation}\index{link adaptation} picks the MCS
from the mobile's channel-quality measurements. And EDGE was the first widely deployed system
with \term{incremental redundancy} (IR), a form of \term{hybrid ARQ}: a block that fails is
not discarded; the receiver keeps its soft values, and the retransmission sends
\emph{different} punctured bits of the same mother code, so the decoder sees a lower-rate code
after each attempt (Chapter~\ref{ch:classiccodes} introduces ARQ). IR makes the choice of
MCS forgiving: a block sent too aggressively costs one more transmission, not a wasted one.
Figure~\ref{fig:ch21_cqi_harq}b, in the LTE section, shows the gain. HSPA, LTE and NR all
inherited IR.

Later releases pushed GSM further. EGPRS2 (Release~7) added 16- and 32-QAM, turbo codes and a
higher symbol rate, and downlink dual-carrier EDGE doubled the peak again, but by then
operators were investing in 3G and LTE.

\begin{worked}[EDGE throughput and what the user saw]
A class-12 EDGE phone receives on four slots. At MCS-9 the peak is $4\times59.2=236.8$\,kb/s.
MCS-9 has a code rate close to 1 and needs a high C/I, roughly 25\,dB or more, so most users
were on MCS-5 to MCS-7: $4\times22.4=90$\,kb/s to $4\times44.8=179$\,kb/s. Subtract RLC/MAC and
IP header overhead, TCP slow start with a 200--300\,ms round trip, and the sharing of the PDCHs
with other users, and a typical user saw something like 100--200\,kb/s. That was enough for
email and simple web pages, and it was the network the first iPhone used in 2007.
\end{worked}

%======================================================================
\section{Third generation: IMT-2000, WCDMA and CDMA2000}
\label{sec:ch21:3g}
%======================================================================

\subsection{The politics of IMT-2000}

By the early 1990s the ITU was working on a single global standard for the next generation,
originally called FPLMTS and renamed IMT-2000, with the idea of one radio interface in one
band everywhere. The World Administrative Radio Conference of 1992 (WARC-92) identified
spectrum around 2\,GHz (1885--2025 and 2110--2200\,MHz), and the requirements were set at
144\,kb/s for vehicular users, 384\,kb/s for pedestrians and 2\,Mb/s indoors. The technology was
expected to be CDMA, whose advantages had been shown by IS-95 and by research in Europe and
Japan.

The single global standard never happened. Europe and Japan wanted a wideband CDMA system
evolved from the GSM core; North American CDMA operators wanted a backward-compatible
evolution of IS-95; China wanted a homegrown TDD system; and patent holders, above all
Qualcomm and Ericsson, fought over the details that would determine whose patents were
essential. In January 1998 ETSI chose \term{wideband CDMA} (WCDMA) for the paired (FDD) bands
and TD-CDMA for unpaired (TDD) bands, under the name UTRA. In December 1998 standards bodies from
Europe, Japan, Korea and the US formed the \term{Third Generation Partnership Project}
(3GPP) to write the specifications for UTRA and the evolved GSM core; a parallel 3GPP2 took on
CDMA2000. A patent dispute between Ericsson and Qualcomm over WCDMA was settled in March 1999,
clearing the way. The ITU eventually blessed a family of five ``IMT-2000'' radio interfaces,
of which WCDMA, CDMA2000 and TD-SCDMA mattered commercially.

The money was even more dramatic than the politics. European governments auctioned 3G
licences in 2000, at the peak of the dot-com bubble: the UK auction raised about
\pounds22.5~billion and the German one about \texteuro50~billion. Operators borrowed heavily
against revenues that were assumed to come from mobile data; when the bubble burst and early
3G handsets turned out to be large, hot and short-lived, the debt crippled the industry for
several years. 3G arrived commercially with NTT DoCoMo's FOMA service in October 2001 and in
Europe in 2002--03, but the data revolution it had been bought for waited for the smartphone.

\begin{history}[How 3GPP works]
3GPP is not a standards body in the legal sense. It is a partnership of seven regional
standards organisations (ARIB and TTC in Japan, ATIS in the US, CCSA in China, ETSI in Europe,
TSDSI in India and TTA in Korea), which transpose its specifications into their own standards.
The work is done by companies---operators, vendors, chipmakers---who send delegates to
technical specification groups (TSG RAN, SA and CT) and their working groups (RAN1 for the
physical layer, RAN2 for the radio protocols, RAN4 for radio performance and so on). Each
meeting processes hundreds or thousands of written \emph{contributions}; decisions are by
consensus, which in practice means persistent negotiation until no company sustains an
objection. Features are packaged into \term{releases}: Release~99 (the first UMTS),
Releases~4--7 (HSPA and its evolution), Release~8 (the first LTE, frozen in December 2008),
Release~10 (LTE-Advanced), Release~15 (the first 5G NR, 2018--19), Release~18 (the first
``5G-Advanced'' release) and so on. Each release goes through study items (technical reports,
TR~3x.9xx) and work items (technical specifications, TS~3x.xxx) and is \emph{frozen} on a
published date, after which only corrections are allowed. The specifications are free to
download from the 3GPP web site, a decision that has done much to spread the knowledge in this
chapter. Every number in this chapter can be checked there, and an engineer who wants to work
in cellular should learn to read them.
\end{history}

\subsection{The WCDMA air interface}

The UMTS Terrestrial Radio Access (UTRA) FDD mode, universally called \term{WCDMA}, uses 5\,MHz
carriers spread at 3.84\,Mchip/s with root-raised-cosine filtering ($\beta=0.22$), so that
$3.84\times1.22=4.7$\,MHz is occupied (one of the problems of Chapter~\ref{ch:baseband} works
this out). Its key
parameters (3GPP TS~25.211--25.214) are:
\begin{itemize}
\item \textbf{Frames and slots.} 10\,ms radio frames of 15 slots, each 2560 chips
(666.7\,$\mu$s). Power control and many physical-layer fields operate per slot.
\item \textbf{Two layers of codes.} \textbf{Channelisation codes}\index{channelisation codes} are orthogonal variable
spreading factor (OVSF) codes, a tree of Walsh-like codes of lengths 4 to 512, which separate
the channels from one transmitter; a high-rate channel takes a short code (and blocks the
longer codes below it in the tree). \textbf{Scrambling codes}\index{scrambling codes} are long Gold sequences that
separate transmitters: on the downlink, each cell is given one of 512 primary scrambling codes
(in 64 groups of eight); on the uplink, each mobile has its own. Chapter~\ref{ch:spreadspectrum}
develops both families.
\item \textbf{Variable rates.} Rates from a few kb/s to 2\,Mb/s are obtained by changing the
spreading factor and the number of codes, with transport channels of 10, 20, 40 or 80\,ms
\emph{transmission time interval} (TTI) multiplexed onto physical channels.
\item \textbf{Coding.} Convolutional codes (constraint length 9) for voice and signalling, and
a rate-1/3 parallel concatenated \term{turbo code} for data (Chapter~\ref{ch:moderncodes}):
the first mass-market use of turbo codes, only six years after their invention.
\item \textbf{Uplink modulation.} The data (DPDCH) and control (DPCCH) channels are placed on
the I and Q branches with different channelisation codes and scrambled with a complex code
designed to reduce the zero crossings and peak-to-average ratio of the handset signal, a scheme
sometimes called HPSK.
\item \textbf{Asynchronous base stations.} Unlike IS-95, base stations need no common time
reference; cells are identified by their scrambling codes, which a phone finds by a three-step
cell search using the primary and secondary synchronisation channels and the common pilot
(Chapter~\ref{ch:sync} works through the equivalent LTE and NR procedures).
\end{itemize}

\subsection{Fast power control and the near--far problem}

In a CDMA uplink all users transmit on the same frequency at the same time, and the base
station separates them by correlation. The processing gain gives each user an advantage of
$10\log_{10}(\mathrm{SF})$\,dB, about 21\,dB for a 12.2\,kb/s voice user on a downlink
channel of spreading factor 128. A user 100\,m from the base has a path loss perhaps 60--70\,dB lower than a
user 5\,km away; without control, the near user's signal would bury the far one completely.
This is the \term{near--far problem}, and its solution is to make every user's signal arrive
at the base with the same power (more precisely, the same SIR target).

Slow (open-loop) control based on the downlink path loss is not enough, because in FDD the
uplink and downlink are 190\,MHz apart and fade independently. WCDMA uses \term{closed-loop
power control} at 1500\,Hz: in every slot the base station measures the received SIR, compares
it with a target, and sends a one-bit transmit-power-control (TPC) command telling the phone to
step up or down by (typically) 1\,dB. An \term{outer loop} in the RNC adjusts the SIR target
slowly to hold the block error rate at its target, typically 1\% for voice. Figure
\ref{fig:ch21_power_control} shows the result in Rayleigh fading. At walking speed the inner
loop follows the fading closely and the received SIR stays within about a decibel of the
target: the phone's transmit power is a mirror image of the fade. At vehicular speeds the
fades are faster than a 1\,dB-per-slot loop can follow, and the SIR distribution approaches
that without power control; there the system relies on the time diversity of interleaving and
coding instead.

\mdcfig{ch21_power_control}{WCDMA closed-loop power control simulated in Rayleigh fading at
2\,GHz, with 1\,dB steps at 1500\,Hz, a one-slot loop delay, a 4\% TPC command error rate and
measurement noise. (a) At 3\,km/h (5.6\,Hz Doppler) the transmit power mirrors the fade and
the received SIR stays near the target (0\,dB). (b) CDF of the received SIR: at 3\,km/h the
loop removes the fading almost completely; at 50 and 120\,km/h (93 and 222\,Hz Doppler) it can
no longer track and the distribution approaches the unfaded-by-control case. The cost at low
speed is transmit power: following deep fades raises the mean transmit power by several
decibels, which raises interference to other cells.}

\begin{keyidea}[In CDMA, power is the shared resource]
In FDMA and TDMA, a user's resource is a channel or a slot, and capacity is the number of
channels. In CDMA every user is interference for every other, so the shared resource is
\emph{received power}, and capacity is reached when the total interference at the base
station (the \term{noise rise}, typically limited to 3--6\,dB above thermal noise) hits its
budget. Power control, voice activity, soft handover and sectorisation all add capacity
because they reduce power. This view---radio resources as an interference budget rather than
a set of channels---carried over directly to the scheduling and interference coordination of
OFDMA systems.
\end{keyidea}

\subsection{Soft handover}

Because neighbouring WCDMA cells use the same frequency, a phone near a cell edge can
receive and transmit to two or three cells at once. The set of cells it is connected to is its
\term{active set}. On the downlink, the phone's rake receiver combines the signals from the
cells as if they were extra multipath components. On the uplink, each cell decodes the phone's
signal and the RNC picks the best copy of each frame (selection combining); when two sectors of
the same site are involved, the base station can combine at the chip level (\emph{softer}
handover). The active set is updated by reporting events: a cell is added when its pilot
$E_c/I_o$ comes within a reporting range (a few decibels) of the best, and removed when it
falls out.

Soft handover gives \term{macrodiversity}: shadowing on paths to different sites is only
partly correlated, so the best of two is much more reliable than either.
Figure~\ref{fig:ch21_soft_handover} shows the effect on the fade margin a cell-edge user needs.
With 8\,dB shadowing, a single cell needs about 10\,dB of margin for 90\% edge coverage; the
best of two cells with shadowing correlation 0.5 needs about 4\,dB less. Because a CDMA
uplink is interference-limited, every decibel of power saved is capacity, and also battery
life. The costs are extra downlink power (two or three cells transmit to the same user),
extra backhaul and RNC processing, and more complicated network planning; typically 20--40\%
of links were in soft handover in a planned WCDMA network.

\mdcfig{ch21_soft_handover}{(a) Macrodiversity with log-normal shadowing ($\sigma=8$\,dB):
probability that a cell-edge user's best link is below the required level, for a single cell
(hard handover) and for the best of two cells with shadowing correlation $\rho$. At 10\%
outage, soft handover saves about 4--6\,dB of fade margin. (b) Pilot $E_c/I_o$ of two cells
along a route between sites 2\,km apart, with each pilot at 10\% of cell power: in the
shaded zone both pilots are within 4\,dB of the best and the phone holds both links.}

\begin{pitfall}[Cell breathing]
In a CDMA network, coverage depends on load. As more users are served, the noise rise at the
base station grows and a cell-edge phone that could just close its link at low load can no
longer reach the required SIR at full power. The cell effectively shrinks, a phenomenon called
\term{cell breathing}. A planner who designs coverage from an empty-network link budget will
find coverage holes in the busy hour. WCDMA link budgets therefore include an
\emph{interference margin} (3\,dB for 50\% load, 6\,dB for 75\%), and admission control rejects
new users before the noise rise exceeds its budget. OFDMA systems with $N=1$ reuse have a
milder version of the same effect: inter-cell interference rises with neighbour load.
\end{pitfall}

\subsection{The UMTS network}

UMTS kept the GSM core network---MSC, HLR, SGSN, GGSN, upgraded---and added a new radio
access network, UTRAN, with the \term{Node~B} (the base station) and the
\term{radio network controller} (RNC) (Figure~\ref{fig:ch21_umts_arch}). The RNC is a much more
powerful node than GSM's BSC. It terminates the radio protocols (RLC, MAC and RRC), performs
soft-handover combining, runs the outer power-control loop and admission control, and talks
to neighbouring RNCs over a new interface, \emph{Iur}, so that soft handover can span RNC
boundaries. The Node~B was deliberately kept simple: modulation, spreading, coding and the
inner power-control loop.

\begin{figure}[htbp]\centering
\begin{tikzpicture}[font=\sffamily\scriptsize,
  box/.style={draw=navy,fill=navy!6,thick,rounded corners=1pt,minimum height=0.8cm,minimum width=1.3cm,align=center},
  db/.style={draw=accent,fill=accent!6,thick,rounded corners=1pt,minimum height=0.7cm,minimum width=1.1cm,align=center},
  pk/.style={draw=practc,fill=practc!8,thick,rounded corners=1pt,minimum height=0.8cm,minimum width=1.3cm,align=center},
  ln/.style={thick,navy}, if/.style={font=\sffamily\tiny,text=accent,midway,above}]
\node[box] (ue) at (0,0) {UE\\(USIM)};
\node[box] (nb1) at (2.3,1.0) {Node B};
\node[box] (nb2) at (2.3,-1.0) {Node B};
\node[box] (rnc1) at (4.8,1.0) {RNC};
\node[box] (rnc2) at (4.8,-1.0) {RNC};
\node[box] (msc) at (7.6,1.0) {MSC/VLR\\(CS domain)};
\node[pk] (sgsn) at (7.6,-1.0) {SGSN\\(PS domain)};
\node[pk] (ggsn) at (10.0,-1.0) {GGSN};
\node[box] (gmsc) at (10.0,1.0) {GMSC};
\node[db] (hlr) at (8.8,0) {HLR/AuC};
\node[draw=gray,thick,ellipse,minimum width=1.3cm] (pstn) at (12.1,1.0) {PSTN};
\node[draw=gray,thick,ellipse,minimum width=1.3cm] (inet) at (12.1,-1.0) {Internet};
\draw[ln,dashed] (ue) -- node[if]{Uu} (nb1); \draw[ln,dashed] (ue) -- (nb2);
\draw[ln] (nb1) -- node[if]{Iub} (rnc1); \draw[ln] (nb2) -- node[if]{Iub} (rnc2);
\draw[ln] (nb2) -- (rnc1);
\draw[ln,accent,thick] (rnc1) -- node[midway,right,font=\sffamily\tiny,text=accent]{Iur} (rnc2);
\draw[ln] (rnc1) -- node[if]{Iu-CS} (msc); \draw[ln,practc] (rnc1) -- (sgsn);
\draw[ln,practc] (rnc2) -- node[if,below]{Iu-PS} (sgsn); \draw[ln] (rnc2) -- (msc);
\draw[ln] (msc) -- (gmsc); \draw[ln] (gmsc) -- (pstn);
\draw[ln,practc] (sgsn) -- node[if]{Gn} (ggsn); \draw[ln,practc] (ggsn) -- node[if]{Gi} (inet);
\draw[ln,accent!70,dotted] (msc) -- (hlr); \draw[ln,accent!70,dotted] (sgsn) -- (hlr);
\node[draw=navy!40,dashed,rounded corners,fit=(nb1)(nb2)(rnc1)(rnc2),inner sep=5pt,label={[font=\sffamily\tiny,text=navy]below:UTRAN: RLC, MAC, RRC, macrodiversity in the RNC}] {};
\end{tikzpicture}
\caption{UMTS Release~99 architecture. The radio network controller terminates the radio
protocols and combines soft-handover links; the Iur interface lets soft handover cross RNC
boundaries. The core keeps GSM's separate circuit-switched (voice) and packet-switched (data)
domains.}
\label{fig:ch21_umts_arch}
\end{figure}

Placing retransmission and scheduling in the RNC was natural for circuit-like dedicated
channels, but it was a poor fit for packet data. Each retransmission crossed the Iub backhaul,
adding tens of milliseconds, and a scheduler in the RNC could not react to fast fading. The
Release~99 phone also had a complicated set of RRC states (CELL\_DCH, CELL\_FACH, CELL\_PCH,
URA\_PCH and idle); moving between them took hundreds of milliseconds to seconds, so early 3G
data sessions felt sluggish even when the peak rate was acceptable. Fixing exactly these
problems was the purpose of HSPA and, later, LTE.

\subsection{CDMA2000 and EV-DO}

The 3GPP2 path evolved IS-95 within its 1.25\,MHz carrier. CDMA2000~1x (first commercial
in South Korea in 2000) doubled voice capacity with better coding and a coherent uplink pilot
and offered packet data at up to 153.6\,kb/s. Its data successor, \term{EV-DO} (Evolution, Data
Optimised, originally Qualcomm's ``high data rate'' system), was a remarkable design and a
forerunner of HSDPA and LTE. It dedicated a whole 1.25\,MHz carrier to data and, on the
downlink, abandoned CDMA's simultaneous code channels: the base station transmits at full
power to one user at a time in 1.67\,ms slots. Each phone reports, every slot, the highest
rate it could currently decode (its data-rate control, DRC), and the base station chooses whom
to serve with a \term{proportional-fair scheduler}, serving each user when its channel is near
its own peak. That exploits \term{multi-user diversity}: with many users fading
independently, there is almost always someone with a good channel
(Chapter~\ref{ch:multipleaccess}). EV-DO Rev.~0 (2002) peaked at 2.4\,Mb/s on the downlink;
Rev.~A (2006) at 3.1\,Mb/s, with a faster uplink and quality-of-service support for VoIP.
3GPP2's planned 4G system, UMB, was abandoned in 2008 when the large CDMA operators, led by
Verizon, chose LTE.

\subsection{TD-SCDMA}

China's \term{TD-SCDMA}, developed by the China Academy of Telecommunications Technology with
Siemens, used 1.6\,MHz TDD carriers at 1.28\,Mchip/s, with 5\,ms subframes, synchronised uplinks,
joint multi-user detection and eight-element ``smart antenna'' arrays at the base station. It
was launched by China Mobile in 2009 and, through it, carried several hundred million users.
Its significance was strategic as much as technical: it built China's capacity in standards and
equipment, which made Chinese companies leading contributors to LTE TDD and then to 5G.

%======================================================================
\section{HSPA: 3G grows up}
\label{sec:ch21:hspa}
%======================================================================

Release~5's \term{High-Speed Downlink Packet Access} (HSDPA, commercial from late 2005)
transformed WCDMA by adopting EV-DO's lessons inside a 5\,MHz carrier. Its features are a
checklist of what packet data needs, and every one of them reappears in LTE.
\begin{itemize}
\item \textbf{A shared channel.} The high-speed downlink shared channel (HS-DSCH) uses up to
fifteen codes of spreading factor~16, shared among users in time and code. There is no power
control on it: the base station transmits with whatever power is left after the other
channels, and adapts the \emph{rate} instead.
\item \textbf{A short TTI.} 2\,ms (three slots) instead of 10--80\,ms, so the scheduler can
react within a few milliseconds.
\item \textbf{Adaptive modulation and coding.} QPSK and 16QAM (64QAM from Release~7), with the
code rate chosen every TTI from a \term{channel quality indicator} (CQI) that the phone reports
as often as every 2\,ms.
\item \textbf{Fast HARQ in the base station.} A new MAC entity in the Node~B (MAC-hs) handles
hybrid ARQ with up to eight parallel stop-and-wait processes, with chase combining or
incremental redundancy, so a retransmission arrives about 12\,ms after the first attempt
instead of crossing to the RNC.
\item \textbf{Scheduling in the base station}, using CQI, buffer state and fairness, usually
some variant of proportional fair.
\item \textbf{No soft handover} on the shared channel: the phone is served by one cell at a
time and switches quickly when another is better.
\end{itemize}

\begin{worked}[HSDPA's peak rate]
With fifteen SF-16 codes, the symbol rate is $15\times3.84\,\mathrm{Mchip/s}/16=3.6$\,Msymbol/s.
With 16QAM (4 bits per symbol) the channel bit rate is $14.4$\,Mb/s. The largest transport block
for a category-10 phone is 27\,952 bits per 2\,ms TTI, 13.98\,Mb/s, which implies an effective
code rate of $27\,952/(2\,\mathrm{ms}\times14.4\,\mathrm{Mb/s})=0.97$: the advertised
``14.4\,Mb/s'' is essentially uncoded and is achievable only at very high SINR, close to the
base station, with no other users. Release~7's 64QAM raises the channel rate to
$3.6\times6=21.6$\,Mb/s (21.1 at the largest transport block), and $2\times2$ MIMO with 64QAM
(Release~8) to 42.2\,Mb/s. Typical early HSDPA users saw 1--3\,Mb/s, ten times the Release~99
rates, and, as important, round-trip times below 100\,ms.
\end{worked}

The uplink followed in Release~6 with \term{HSUPA} (the enhanced dedicated channel, E-DCH):
2 or 10\,ms TTIs, HARQ, and a Node-B scheduler that controls each phone's permitted power ratio
through absolute and relative grants. HSUPA keeps power control and soft handover, since the
uplink remains interference-limited, and peaks at 5.76\,Mb/s (11.5\,Mb/s with 16QAM in
Release~7). \term{HSPA+} (Releases~7--11) added 64QAM, $2\times2$ MIMO, \emph{continuous packet
connectivity} (discontinuous transmission and reception, so smartphones could stay connected
without draining their batteries), a flatter architecture that can bypass the SGSN or move RNC
functions into the Node~B, and multi-carrier operation: dual-carrier HSDPA (42\,Mb/s,
Release~8), dual-carrier with MIMO (84\,Mb/s, Release~9) and four carriers (168\,Mb/s,
Release~10). By then LTE was deployed, and HSPA+ became the fallback layer.

\begin{history}[The iPhone and the data tsunami]
Apple's first iPhone (June 2007) supported only GSM/EDGE. It was not the first smartphone, but
it was the first that people used the way they used a computer: a full web browser,
maps, YouTube, email with attachments. The iPhone~3G (July 2008) added HSDPA, and Apple's App
Store opened the same month; Android phones followed from late 2008. The effect on networks was
immediate. AT\&T, the exclusive US iPhone carrier, saw its mobile data traffic grow by roughly
fifty times in three years, and its networks in New York and San Francisco became a byword
for dropped calls and stalled downloads in 2009. The problem was not only volume. Smartphone
apps sent frequent small keep-alive messages, each of which moved the phone from idle to a
connected RRC state and back, so the signalling load per megabyte was far higher than for a
laptop modem; networks experienced \emph{signalling storms}, and phones invented ``fast
dormancy'' to drop back to idle early, which made it worse. The engineering response took
several forms: HSPA+ with continuous packet connectivity, LTE with its simpler two-state RRC and
fast idle-to-connected transitions, and, a decade later, NR's RRC\_INACTIVE state. The
commercial response was the end of unlimited data plans in many markets. The smartphone, not
3G, created the mobile data business, and 4G was designed for it.
\end{history}

%======================================================================
\section{Fourth generation: LTE}
\label{sec:ch21:lte}
%======================================================================

\subsection{Design goals and the decisions of 2004--2008}

3GPP began work on the ``Long Term Evolution'' of UMTS at a workshop in Toronto in November
2004, under pressure from two directions: WiMAX (IEEE~802.16e), an OFDMA system from the
computer industry that promised high rates with a simple IP architecture, and the realisation
that HSPA's CDMA core would not scale to 20\,MHz carriers. The targets, recorded in
3GPP TR~25.913, were aggressive for their time:
\begin{itemize}
\item peak rates of 100\,Mb/s downlink and 50\,Mb/s uplink in 20\,MHz (5 and 2.5\,b/s/Hz);
\item average spectral efficiency three to four times that of HSDPA Release~6 on the
downlink and two to three times HSUPA on the uplink;
\item user-plane latency below 5\,ms one way across the radio network, and an idle-to-active
transition below 100\,ms;
\item scalable bandwidth from 1.4 to 20\,MHz, paired (FDD) and unpaired (TDD) spectrum;
\item mobility up to 350\,km/h, optimised for 0--15\,km/h;
\item a packet-only, simplified architecture with fewer nodes, and lower cost per bit.
\end{itemize}
Release~8, the first LTE release, was frozen in December 2008. The first commercial networks
opened in Stockholm and Oslo in December 2009 (TeliaSonera), and the US operators Verizon and
MetroPCS followed in 2010. By 2012 LTE was the fastest-deployed mobile technology in history,
and it became the first truly global standard: the CDMA2000 operators, the WCDMA operators,
China's TD-SCDMA operator and the WiMAX hopefuls all converged on it.

\subsection{Why OFDMA down and SC-FDMA up}

The single most important decision was to abandon CDMA for \term{OFDMA} on the downlink.
Chapter~\ref{ch:ofdm} develops OFDM in detail; the reasons it won are worth stating in
cellular terms.
\begin{enumerate}
\item \textbf{Wide bandwidths without an equaliser.} At 20\,MHz, a delay spread of
5\,$\mu$s spans about 150 chips; a rake or chip equaliser over that many taps per antenna is
expensive, and its performance degrades with spatial multiplexing. OFDM with a cyclic prefix
turns the channel into parallel flat subchannels, each equalised with one complex multiply.
\item \textbf{Scheduling in frequency as well as time.} In a frequency-selective channel, each
user's good subcarriers differ. An OFDMA scheduler can give each user the resource blocks
where its channel is strong---multi-user diversity in two dimensions.
\item \textbf{MIMO.} Spatial multiplexing is simple per subcarrier: a small matrix per resource
element rather than a space--time equaliser (Chapter~\ref{ch:mimo}).
\item \textbf{Scalability.} The same subcarrier spacing and symbol length work for every
bandwidth from 1.4 to 20\,MHz; only the number of subcarriers and the FFT size change.
\end{enumerate}
OFDM's weakness is its high peak-to-average power ratio, which matters most in the handset
power amplifier. The uplink therefore uses \term{SC-FDMA}, the DFT-spread OFDM of
Section~\ref{sec:ch17:dfts}: each user's symbols are DFT-precoded before mapping to a
contiguous block of subcarriers, so the transmitted signal is a single-carrier waveform with
a cyclic prefix. Its PAPR is several decibels lower than OFDM's for QPSK (Lab~11, Section~5),
which translates directly into handset power-amplifier efficiency, battery life and uplink
coverage, while the base station keeps OFDMA's frequency-domain equalisation and scheduling.

\begin{keyidea}[The radio resource is a time--frequency grid]
In LTE and NR the scheduler's job is to fill a grid. Its unit, the \term{resource block}
(RB), is 12 subcarriers (180\,kHz in LTE) by one slot; its granularity in time is the
1\,ms subframe (LTE) or the slot or mini-slot (NR). Every millisecond, the base station
decides which user gets which RBs, with which modulation, code rate and MIMO precoding, and
tells them on a control channel. All of the system's intelligence---link adaptation,
HARQ, MIMO, quality of service, interference coordination---lives in that decision, which
is why scheduler design is where vendors compete and where much of the real-world
performance difference between networks comes from.
\end{keyidea}

\subsection{A flat, all-IP architecture}

LTE's system architecture (the \term{Evolved Packet System}, EPS) removed the RNC and the
circuit-switched domain (Figure~\ref{fig:ch21_lte_arch}).
\begin{itemize}
\item The radio network, E-UTRAN, has one kind of node, the \term{eNodeB} (eNB), which hosts
the whole radio protocol stack---PHY, MAC, RLC, PDCP and RRC---and makes all radio
resource decisions: scheduling, HARQ, link adaptation, admission and handover. Neighbouring
eNBs talk directly over the X2 interface for handover and interference coordination.
\item The \term{Evolved Packet Core} (EPC) separates control from user plane. The
\term{mobility management entity} (MME) handles signalling: attach, authentication (with the
\term{home subscriber server}, HSS), tracking-area management, paging and bearer set-up. The
\term{serving gateway} (S-GW) anchors the user plane during handovers between eNBs, and the
\term{PDN gateway} (P-GW) allocates the user's IP address and connects to external networks;
the policy and charging rules function (PCRF) sets quality-of-service rules.
\item User data travels in \term{bearers}, tunnels with a defined quality of service (a QoS
class identifier, QCI), from the phone through the eNB (in GTP-U tunnels on S1-U) to the
S-GW and P-GW.
\end{itemize}
Putting everything radio-related in the eNB cut latency: HARQ retransmissions happen at the
radio site, 8\,ms after the first attempt, and the scheduler sees the channel a few
milliseconds old. Handover became ``hard'' (break-before-make) but fast, with data forwarded
from the old eNB to the new one over X2 to avoid loss.

\begin{figure}[htbp]\centering
\begin{tikzpicture}[font=\sffamily\scriptsize,
  box/.style={draw=navy,fill=navy!6,thick,rounded corners=1pt,minimum height=0.8cm,minimum width=1.25cm,align=center},
  db/.style={draw=accent,fill=accent!6,thick,rounded corners=1pt,minimum height=0.7cm,minimum width=1.1cm,align=center},
  up/.style={draw=practc,fill=practc!8,thick,rounded corners=1pt,minimum height=0.8cm,minimum width=1.25cm,align=center},
  cp/.style={thick,accent!80,dashed}, upl/.style={thick,practc},
  if/.style={font=\sffamily\tiny,text=black!70,midway,above}]
\node[box] (ue) at (0,0) {UE};
\node[box] (enb1) at (2.3,0.9) {eNB};
\node[box] (enb2) at (2.3,-0.9) {eNB};
\node[db] (mme) at (5.2,1.6) {MME};
\node[up] (sgw) at (5.2,-0.6) {S-GW};
\node[up] (pgw) at (7.7,-0.6) {P-GW};
\node[db] (hss) at (7.7,1.6) {HSS};
\node[db] (pcrf) at (10.0,0.6) {PCRF};
\node[draw=gray,thick,ellipse,minimum width=1.6cm,align=center] (pdn) at (10.2,-0.9) {Internet,\\IMS};
\draw[thick,navy,dashed] (ue) -- node[if]{LTE-Uu} (enb1); \draw[thick,navy,dashed] (ue) -- (enb2);
\draw[thick,navy] (enb1) -- node[midway,left,font=\sffamily\tiny]{X2} (enb2);
\draw[cp] (enb1) -- node[if,sloped]{S1-MME} (mme); \draw[cp] (enb2) -- (mme);
\draw[upl] (enb1) -- node[if,sloped,pos=0.6]{S1-U} (sgw); \draw[upl] (enb2) -- (sgw);
\draw[cp] (mme) -- node[midway,left,font=\sffamily\tiny]{S11} (sgw);
\draw[cp] (mme) -- node[if]{S6a} (hss);
\draw[upl] (sgw) -- node[if]{S5/S8} (pgw);
\draw[upl] (pgw) -- node[if,below]{SGi} (pdn);
\draw[cp] (pgw) -- node[midway,above,font=\sffamily\tiny,sloped]{Gx} (pcrf);
\node[draw=navy!40,dashed,rounded corners,fit=(enb1)(enb2),inner sep=5pt,label={[font=\sffamily\tiny,text=navy]below:E-UTRAN}] {};
\node[draw=navy!40,dashed,rounded corners,fit=(mme)(sgw)(pgw)(hss)(pcrf),inner sep=5pt,label={[font=\sffamily\tiny,text=navy]below:Evolved Packet Core}] {};
\node[font=\sffamily\tiny,text=accent,align=left] at (0.3,1.8) {- - - control plane};
\node[font=\sffamily\tiny,text=practc,align=left] at (0.3,1.5) {------ user plane};
\end{tikzpicture}
\caption{The LTE/EPS architecture (3GPP TS 23.401). One radio node, the eNB, holds the whole
radio stack; the core separates signalling (MME) from the user-plane gateways (S-GW, P-GW).
There is no circuit-switched domain: voice is an IP application (VoLTE) on the IMS.}
\label{fig:ch21_lte_arch}
\end{figure}

\subsection{The radio protocol stack}

Figure~\ref{fig:ch21_lte_stack} shows the radio protocols, which NR kept with small changes.
From the top:
\begin{itemize}
\item \textbf{RRC} (radio resource control, control plane only) configures everything below
it: bearers, measurements, handovers, MIMO modes, DRX cycles. It defines the phone's
\term{RRC states}: RRC\_IDLE, in which the phone is known only to the core at the level of a
tracking area, listens for paging and chooses its own cell; and RRC\_CONNECTED, in which the
eNB knows it, schedules it and controls its handovers. LTE's two states replaced UMTS's five.
\item \textbf{NAS} (non-access stratum) runs between the phone and the MME, transparent to
the eNB: attach, authentication, security mode, tracking-area updates and bearer management.
\item \textbf{PDCP} (packet data convergence protocol) performs header compression (robust
header compression, ROHC, which shrinks a 40- or 60-byte IP/UDP/RTP header to a few bytes),
ciphering, integrity protection of signalling, and in-order delivery and duplicate removal
at handover.
\item \textbf{RLC} (radio link control) segments and reassembles packets to fit the transport
block chosen by the scheduler, and in acknowledged mode runs an ARQ that catches the residual
errors of HARQ.
\item \textbf{MAC} multiplexes logical channels into transport blocks, runs HARQ, the random
access procedure and timing advance, and---in the eNB---the scheduler.
\item \textbf{PHY} handles coding, modulation, MIMO, OFDM and the physical channels described
below.
\end{itemize}

\begin{figure}[htbp]\centering
\begin{tikzpicture}[font=\sffamily\scriptsize,
  l/.style={draw=navy,fill=navy!6,thick,minimum width=1.55cm,minimum height=0.42cm,align=center},
  c/.style={draw=accent,fill=accent!6,thick,minimum width=1.55cm,minimum height=0.42cm,align=center},
  u/.style={draw=practc,fill=practc!8,thick,minimum width=1.55cm,minimum height=0.42cm,align=center}]
% control plane
\node[font=\sffamily\bfseries\scriptsize,text=navy] at (2.6,3.25) {Control plane};
\foreach \x/\nm in {0/UE,1.75/eNB,3.5/MME}{\node[font=\sffamily\bfseries\tiny] at (\x+0.8,2.85) {\nm};}
\node[c] at (0.8,2.45) {NAS}; \node[c] at (4.3,2.45) {NAS};
\node[c] at (0.8,2.0) {RRC}; \node[c] at (2.55,2.0) {RRC};
\foreach \y/\nm in {1.55/PDCP,1.1/RLC,0.65/MAC,0.2/PHY}{\node[l] at (0.8,\y) {\nm}; \node[l] at (2.55,\y) {\nm};}
\node[l] at (4.3,2.0) {S1-AP}; \node[l] at (4.3,1.55) {SCTP}; \node[l] at (4.3,1.1) {IP}; \node[l] at (4.3,0.65) {L2/L1};
\draw[dotted,thick,accent] (1.6,2.45) -- (3.5,2.45);
% user plane
\node[font=\sffamily\bfseries\scriptsize,text=navy] at (9.2,3.25) {User plane};
\foreach \x/\nm in {6.6/UE,8.35/eNB,10.1/S-GW / P-GW}{\node[font=\sffamily\bfseries\tiny] at (\x+0.8,2.85) {\nm};}
\node[u] at (7.4,2.45) {IP (app)}; \node[u] at (10.9,2.45) {IP};
\foreach \y/\nm in {1.55/PDCP,1.1/RLC,0.65/MAC,0.2/PHY}{\node[l] at (7.4,\y) {\nm}; \node[l] at (9.15,\y) {\nm};}
\node[l] at (9.15,2.0) {GTP-U}; \node[l] at (10.9,2.0) {GTP-U}; \node[l] at (10.9,1.55) {UDP/IP}; \node[l] at (10.9,1.1) {L2/L1};
\draw[dotted,thick,practc] (8.2,2.45) -- (10.1,2.45);
\draw[navy!50,thick] (5.55,-0.1) -- (5.55,3.3);
\node[font=\sffamily\tiny,text=black!60,align=center] at (2.6,-0.35) {LTE-Uu \hspace{1.0cm} S1-MME};
\node[font=\sffamily\tiny,text=black!60,align=center] at (9.2,-0.35) {LTE-Uu \hspace{1.2cm} S1-U};
\end{tikzpicture}
\caption{LTE radio protocol stacks. In the control plane, RRC terminates in the eNB and NAS in
the MME; in the user plane, PDCP, RLC and MAC terminate in the eNB and user packets are
tunnelled to the gateways in GTP-U. NR uses the same layering with a new SDAP layer above PDCP
for QoS-flow mapping.}
\label{fig:ch21_lte_stack}
\end{figure}

\subsection{Frame structure and the resource grid}

LTE's numerology is defined in 3GPP TS~36.211 in units of
$T_s=1/(15\,000\times2048)\,\mathrm{s}\approx32.6$\,ns, the sample period of a 2048-point FFT at
30.72\,MS/s. The \term{radio frame} is 10\,ms, divided into ten 1\,ms \term{subframes}, each of
two 0.5\,ms \term{slots}. With the normal cyclic prefix a slot holds seven OFDM symbols (the
first with a 160-sample, 5.2\,$\mu$s prefix, the others 144 samples, 4.7\,$\mu$s); with the
extended prefix (16.7\,$\mu$s) it holds six. Subcarriers are 15\,kHz apart. A resource block is
12 subcarriers by one slot, and the scheduler allocates pairs of RBs, one subframe long. The
channel bandwidths of 1.4, 3, 5, 10, 15 and 20\,MHz contain 6, 15, 25, 50, 75 and 100 RBs,
about 90\% occupancy for the wider ones. Chapter~\ref{ch:ofdm} works through the numerology's
design (Section~\ref{sec:ch17:design}) and the resource grid (Section~\ref{sec:ch17:grid}). In
TDD (frame structure type~2), the subframes are assigned to downlink or uplink by one of seven
uplink--downlink configurations, with a \emph{special subframe} (downlink part, guard period,
uplink part) at each downlink-to-uplink switch.

\mdcfig[0.92]{ch21_lte_grid}{Top: the LTE FDD radio frame, with the synchronisation signals in
subframes 0 and 5 and the broadcast channel in subframe 0. Bottom: subframe 0 of a 1.4\,MHz
carrier with two cell-specific reference-signal (CRS) ports and a three-symbol control region,
following TS 36.211. The PSS and SSS occupy the 62 subcarriers around the carrier centre in the
last two symbols of the first slot; the PBCH occupies the first four symbols of the second
slot over the central 72 subcarriers, skipping positions reserved for the CRS of four ports.
The CRS are present in every subframe, whether or not anyone is being served.}

\subsection{Physical channels and reference signals}

Figure~\ref{fig:ch21_lte_grid} shows how the downlink grid is populated. The physical signals
and channels of LTE (TS~36.211) are as follows.
\begin{description}[style=nextline,leftmargin=1.2em,font=\sffamily\bfseries\small\color{navy}]
\item[Synchronisation signals.] The \term{primary synchronisation signal} (PSS), one of three
length-63 Zadoff--Chu sequences, and the \term{secondary synchronisation signal} (SSS), an
interleaved pair of m-sequences, are sent twice per frame on the central 62 subcarriers. Together
they give the 504 physical cell identities, slot and frame timing and the cyclic-prefix length.
Chapter~\ref{ch:sync} (Section~\ref{sec:ch10:frame}) analyses the search.
\item[Cell-specific reference signals (CRS).] Pilots on a diagonal pattern, for one, two or four
antenna ports, in every subframe across the whole bandwidth. Phones use them for channel
estimation, demodulation, CSI measurement and mobility measurements. They are LTE's most
characteristic feature and, as the NR section explains, its biggest regret: they are
transmitted whether or not there is traffic, which costs energy and interference and cannot be
beamformed.
\item[PBCH.] The physical broadcast channel carries the 24-bit \emph{master information
block} (bandwidth, PHICH configuration, frame number) with a 40\,ms TTI, repeated in four
self-decodable parts. Everything else about the cell is sent in \emph{system information
blocks} on the PDSCH.
\item[Control region.] The first one to three OFDM symbols of each subframe (two to four for
1.4\,MHz) carry the PCFICH (how many control symbols), the PHICH (ACK/NACK for uplink HARQ)
and the \term{PDCCH}, which carries \emph{downlink control information} (DCI): the scheduling
assignments and grants. A PDCCH occupies 1, 2, 4 or 8 \emph{control channel elements} of 36
resource elements each, depending on the user's channel quality; its CRC is scrambled with the
user's identifier, and each phone blindly decodes up to several tens of candidate positions
every subframe to find its own messages.
\item[PDSCH.] The shared data channel: everything else, including paging and system
information.
\item[Later additions.] UE-specific demodulation reference signals for beamforming (Release~9),
channel-state-information reference signals (CSI-RS) for up to eight ports (Release~10),
positioning reference signals, and the enhanced PDCCH, which moved control into the data
region (Release~11).
\end{description}
On the uplink, the \term{PUCCH} carries uplink control (HARQ ACK/NACK, scheduling requests,
CQI and precoding reports) on resource blocks at the two edges of the band, hopping between
them at the slot boundary for frequency diversity. The \term{PUSCH} carries data, with a
demodulation reference signal in the middle symbol of each slot (another midamble, as in GSM).
The \term{PRACH} carries random-access preambles, 64 per cell, built from length-839 Zadoff--Chu
sequences. \textbf{Sounding reference signals}\index{sounding reference signals} (SRS) let the eNB measure the uplink channel
across the band for frequency-selective scheduling and, in TDD, for downlink beamforming by
reciprocity.

\subsection{Coding, HARQ and link adaptation}

Data channels use the rate-1/3 \term{turbo code} of UMTS with a new
\emph{quadratic permutation polynomial} (QPP) interleaver, designed for parallel decoding
(Chapter~\ref{ch:moderncodes}), and 188 block sizes from 40 to 6144 bits. A transport block
gets a 24-bit CRC, is segmented into code blocks (each with its own 24-bit CRC so that a decoder
can stop early) and passes through \term{circular-buffer rate matching}: the systematic and
parity bits are written into a buffer, and each transmission reads a contiguous window starting
at one of four \emph{redundancy versions}. The first transmission (RV0) is mostly systematic
bits; later RVs add parity, implementing incremental redundancy. Control channels use a
tail-biting convolutional code (Chapter~\ref{ch:classiccodes}).

\textbf{Hybrid ARQ}\index{hybrid ARQ} in LTE FDD uses eight parallel stop-and-wait processes with a fixed timing:
data in subframe $n$, ACK/NACK in subframe $n+4$, retransmission at the earliest in $n+8$. The
4\,ms allows about 3\,ms of decoding time in the phone. The downlink HARQ is asynchronous and
adaptive (the eNB may retransmit any time, with a different allocation and modulation); the
uplink is synchronous (the retransmission happens exactly eight subframes later unless
cancelled), which saves signalling.

\textbf{Link adaptation}\index{link adaptation} closes the loop. The phone reports a 4-bit \term{CQI}: the index of the
highest-rate entry of TS~36.213 Table~7.2.3-1 that it could receive with a block error rate
of at most 10\%, together with a precoding matrix indicator (PMI) and rank indicator (RI) for
MIMO. The eNB maps the report to one of 29 modulation and coding schemes and a transport block
size, usually adding an outer-loop offset that nudges the actual first-transmission BLER to the
10\% target using the HARQ ACK statistics. Why aim for 10\% and not 0.1\%? Because HARQ makes
errors cheap. Figure~\ref{fig:ch21_cqi_harq} shows both halves of this design: the CQI table
tracks Shannon's curve a few decibels to the right, and incremental redundancy makes the
throughput degrade gracefully when the first transmission was too optimistic.

\mdcfig{ch21_cqi_harq}{(a) The 15 entries of the LTE CQI table (TS 36.213 Table 7.2.3-1),
QPSK to 64QAM with code rates from 0.076 to 0.93, placed at the SINR at which a code a fixed
2\,dB from capacity would decode (a common system-simulation abstraction), and the resulting
throughput staircase at 10\% BLER. (b) Why HARQ matters: a transmission sent at 3\,b/s/Hz fails
below about 10\,dB SNR, but with up to four transmissions, chase combining (which adds SNR)
and especially incremental redundancy (which adds parity, and hence mutual information) recover
a large fraction of the rate the channel could support.}

\begin{worked}[LTE's peak rate: 20\,MHz, $4\times4$, 64QAM]
A 20\,MHz carrier has 100 RBs, $100\times12=1200$ subcarriers, and with the normal prefix
14 symbols per 1\,ms subframe: $1200\times14=16\,800$ resource elements per millisecond.

\emph{Overheads.} Take the smallest control region, one symbol: 1200 REs. With four CRS ports,
each RB pair carries 24 CRS REs, of which 4 fall in the control symbol, so 20 of them in the data
region: $100\times20=2000$ REs. Synchronisation and broadcast channels occupy subframes 0 and 5
only and are about 1\% on average; ignore them. The PDSCH has about $16\,800-1200-2000=13\,600$
REs per millisecond, an overhead of 19\%.

\emph{Bits.} With 64QAM (6 bits) and four spatial layers, $13\,600\times6\times4=326\,400$ coded
bits per millisecond, 326\,Mb/s. At the highest code rate (about 0.92) that is about 300\,Mb/s.
The largest transport block for a category-5 phone is 299\,552 bits per subframe: 299.6\,Mb/s.

\emph{Reality.} Few phones had four receive antennas, and a peak rate needs the whole cell,
SINR above about 25\,dB and a rank-4 channel. Category-3 and -4 phones ($2\times2$) peaked at
100 and 150\,Mb/s, and typical user rates in a loaded network were 10--50\,Mb/s. The average
spectral efficiency of a real LTE macro network is of the order of 1.5--2.5\,b/s/Hz per cell,
roughly a tenth of the peak.
\end{worked}

\subsection{MIMO transmission modes}

LTE was the first cellular standard designed around multiple antennas
(Chapter~\ref{ch:mimo}). Its \term{transmission modes} (TM) define how the PDSCH uses them;
Table~\ref{tab:ch21:tm} summarises them. The early modes (TM2--TM4) rely on the CRS and a
codebook of precoders known to both sides; the later ones (TM7--TM10) use UE-specific
demodulation reference signals, so that the eNB can apply any precoder---including beams
computed from uplink sounding in TDD---without telling the phone what it is.

\begin{table}[htbp]\centering\small
\caption{LTE downlink transmission modes (TS 36.213).}
\label{tab:ch21:tm}
\begin{tabularx}{\linewidth}{@{}l l X@{}}
\toprule
Mode & Release & Description\\
\midrule
TM1 & 8 & Single antenna port\\
TM2 & 8 & Transmit diversity: space--frequency block code (Alamouti on pairs of subcarriers)\\
TM3 & 8 & Open-loop spatial multiplexing with large-delay cyclic delay diversity (high speed)\\
TM4 & 8 & Closed-loop spatial multiplexing with codebook precoding (PMI/RI feedback)\\
TM5 & 8 & Multi-user MIMO with codebook precoding\\
TM6 & 8 & Closed-loop rank-1 precoding (beamforming with the codebook)\\
TM7 & 8 & Single-layer beamforming with UE-specific reference signal (port 5), mainly TDD\\
TM8 & 9 & Dual-layer beamforming (ports 7 and 8), SU/MU-MIMO\\
TM9 & 10 & Up to 8 layers with CSI-RS for measurement and DM-RS for demodulation\\
TM10 & 11 & As TM9, with coordinated multipoint (CoMP) and quasi-co-location signalling\\
\bottomrule
\end{tabularx}
\end{table}

\subsection{LTE-Advanced and LTE-Advanced Pro}

The ITU's requirements for IMT-Advanced (ITU-R M.2134, 2008)---including 1\,Gb/s peak for
low mobility, peak spectral efficiencies of 15\,b/s/Hz downlink and 6.75\,b/s/Hz uplink and
bandwidths up to 40\,MHz---were beyond Release~8. \term{LTE-Advanced} (Release~10, 2011) met
them and was accepted, alongside IEEE~802.16m, as an IMT-Advanced technology: the first
``real'' 4G by the ITU's definition. Its main additions, and those of the following releases,
were:
\begin{itemize}
\item \textbf{Carrier aggregation} (CA): a phone combines up to five \emph{component carriers}
of up to 20\,MHz each, contiguous or in different bands, into one connection of up to
100\,MHz; Release~13 raised the limit to 32 carriers. Each carrier keeps its Release-8
structure, so legacy phones can use any of them. CA was LTE-Advanced's commercial success,
because operators' spectrum is fragmented across many bands. A primary cell (PCell) carries
the control signalling and secondary cells (SCells) add capacity.
\item \textbf{More antennas}: up to eight layers on the downlink and four on the uplink, and in
Release~13 ``full-dimension'' MIMO with up to 16 (later 32) CSI-RS ports for two-dimensional
arrays that steer in elevation as well as azimuth.
\item \textbf{Heterogeneous networks}: small cells under the macro layer, with
\term{enhanced inter-cell interference coordination} (eICIC). The macro cell transmits
\emph{almost blank subframes}, in which it sends only reference signals, and pico cells
schedule their cell-edge users (pushed outwards by a \emph{cell range expansion} bias) in
those subframes.
\item \textbf{Relays}: decode-and-forward relay nodes with in-band wireless backhaul.
\item \textbf{Coordinated multipoint} (CoMP, Release~11): coordinated scheduling and beamforming
across sites, or joint transmission from several. CoMP needed fast, low-latency backhaul and
proved hard to make pay in practice; its ideas returned in NR's multi-TRP transmission.
\item \textbf{Release~12--14}: 256QAM, dual connectivity (simultaneous connection to two eNBs,
the basis of 5G non-standalone operation), device-to-device communication for public safety,
licensed-assisted access (LAA) to the 5\,GHz unlicensed band with listen-before-talk, and
vehicle-to-everything (V2X) sidelinks. Releases~13 and~14 were branded
\term{LTE-Advanced Pro}.
\item \textbf{IoT variants} (Release~13): \term{LTE-M} (eMTC, 1.4\,MHz, about 1\,Mb/s) and
\term{NB-IoT} (180\,kHz, tens of kb/s, a maximum coupling loss of 164\,dB, about 20\,dB
better than GSM) for low-cost, long-battery-life devices; Chapter~\ref{ch:wifi} compares them
with LoRa and other IoT radios.
\end{itemize}
``Gigabit LTE'' phones (category~16 and above, from 2017) combined carrier aggregation of three
to five carriers, $4\times4$ MIMO on some of them and 256QAM to reach a peak of about 1\,Gb/s.

\subsection{Voice over LTE and the IMS}

LTE has no circuit-switched domain, so voice had to become an IP application. The first LTE
phones used \term{circuit-switched fallback} (CSFB): on an incoming or outgoing call, the phone
dropped to 2G or 3G for the duration, at the cost of a few seconds of call set-up time and no
LTE data during the call. The target solution, \term{VoLTE} (specified in the GSMA's IR.92
profile and first launched in 2012), carries voice as RTP packets over a dedicated bearer with
QCI~1 (guaranteed bit rate, a 100\,ms packet delay budget) and signalling over the
\term{IP Multimedia Subsystem} (IMS), a SIP-based core of call-session control functions.
VoLTE needed every layer of the stack to cooperate:
\begin{itemize}
\item AMR-WB or, from 2014, the EVS codec (Chapter~\ref{ch:sourcecoding}), sending a packet every
20\,ms;
\item ROHC header compression, without which the 40 to 60-byte IP/UDP/RTP headers would be
larger than the 30-odd-byte speech payload;
\item \term{semi-persistent scheduling} (SPS), which allocates a fixed resource every 20\,ms
with one control message instead of one per packet;
\item \emph{TTI bundling}, which sends each uplink packet in four consecutive subframes to give
cell-edge users enough energy per packet;
\item DRX, so the phone's receiver sleeps between packets;
\item single radio voice call continuity (SRVCC), which hands a VoLTE call over to a 2G or 3G
circuit when the phone leaves LTE coverage.
\end{itemize}

\begin{worked}[How many VoLTE calls fit in 10\,MHz?]
Take a 10\,MHz FDD carrier (50 RBs) and AMR-WB at 12.65\,kb/s: 253 bits every 20\,ms. Adding the
RTP payload header, a compressed ROHC header of about 3 bytes, PDCP, RLC and MAC headers and the
24-bit CRC gives roughly 40 bytes, 320 bits, per transport block.

\emph{Resource-limited capacity.} Of the 168 REs in an RB pair, the three-symbol control region
and the CRS of two ports take about 48, leaving about 120. At a typical mid-cell modulation and
coding (16QAM at rate about 0.4, 1.6 bits per RE) one RB pair carries about 190 bits, so a voice
packet needs two RB pairs and 50 RBs carry about 25 packets per subframe. Each call needs one
packet every 20 subframes, so $25\times20=500$ calls could be served if every call talked
continuously. With a voice-activity factor of about 0.5 and small silence-descriptor packets,
the limit approaches 800--900 calls.

\emph{Control-limited capacity.} With dynamic scheduling, each packet also needs a DCI on the
PDCCH. A three-symbol control region in 10\,MHz holds about 40 control channel elements; at an
average aggregation of two and with uplink grants sharing the space, only about 10 downlink
assignments fit per subframe, so the cell saturates at about $10\times20=200$ simultaneously
talking calls, perhaps 400 calls with voice activity. That is why VoLTE uses semi-persistent
scheduling, which removes most of the per-packet control.

\emph{Comparison.} The same 10\,MHz as GSM, with 4/12 reuse, gives each sector about four
carriers: some 30 full-rate traffic channels, or 60 with half-rate AMR. These are rough
estimates---system simulations with realistic SINR distributions, packet-loss targets and
uplink limits give lower numbers for VoLTE---but an order-of-magnitude gain in voice capacity per
megahertz over GSM is the right picture.
\end{worked}

\begin{inpractice}[Reading an LTE cell from your phone]
Engineering screens and apps that display a phone's modem state (on Android, several free apps
read the radio interface layer) show LTE's parameters live: the EARFCN (channel number) and
band, the physical cell identity (0--503), the bandwidth, and the measurements RSRP (reference
signal received power, the average power of one CRS resource element, typically $-80$\,dBm
near a site and $-120$\,dBm at the edge of coverage), RSRQ (RSRP relative to total received
power, from about $-3$ to $-20$\,dB) and SINR. RSRP is a per-resource-element quantity, so it is
about $10\log_{10}(1200)\approx31$\,dB below the total received power of a fully loaded 20\,MHz
carrier: an RSRP of $-100$\,dBm is a respectable signal, not a weak one. With a USRP B200 and an
open-source LTE receiver you can decode the MIB and system information of a nearby cell, as
suggested in Lab~11.
\end{inpractice}

%======================================================================
\section{Fifth generation: 5G NR}
\label{sec:ch21:nr}
%======================================================================

\subsection{IMT-2020: three kinds of user}

5G's starting point was not a bottleneck in LTE's data rate, which by 2015 was adequate for
most smartphone use, but a broader ambition. The ITU's vision for IMT-2020 (Recommendation
ITU-R M.2083, 2015) defined three usage scenarios, often drawn as the corners of a triangle:
\begin{itemize}
\item \term{enhanced mobile broadband} (eMBB): more of what LTE did, with peak rates of
gigabits per second, much higher capacity per square kilometre and a more uniform user
experience;
\item \term{ultra-reliable low-latency communication} (URLLC): industrial control, vehicle
coordination and remote operation, with air-interface latency of about 1\,ms and reliability
of $1-10^{-5}$ for a small packet;
\item \term{massive machine-type communication} (mMTC): very large numbers of cheap, low-power
devices, up to a million per square kilometre.
\end{itemize}
Figure~\ref{fig:ch21_imt_capabilities} compares the eight ``key capabilities'' of M.2083 with
those of IMT-Advanced. The minimum technical requirements that a candidate had to meet
(ITU-R M.2410) include a peak rate of 20\,Gb/s downlink and 10\,Gb/s uplink, peak spectral
efficiencies of 30 and 15\,b/s/Hz, a user-plane latency of 4\,ms for eMBB and 1\,ms for URLLC,
a control-plane latency of 20\,ms, mobility up to 500\,km/h and support for at least 100\,MHz of
bandwidth, up to 1\,GHz above 6\,GHz. 3GPP's NR, submitted with LTE as a combined proposal, was
approved as IMT-2020 in 2020 (ITU-R M.2150).

\mdcfig{ch21_imt_capabilities}{The key capabilities of IMT-2020 relative to IMT-Advanced, from
ITU-R M.2083. Note the shape of the target: the peak rate grows by 20, but the biggest factors
are in area traffic capacity, connection density and network energy efficiency, which call for
dense small cells, wide bandwidths and lean signalling rather than better modulation.}

\subsection{Design principles of NR}

The 3GPP study for NR ran from 2015 to 2017, and the first specification, Release~15, was
completed in three steps: a non-standalone version in December 2017, the standalone version in
June 2018 and a ``late drop'' in March 2019. Commercial service started in South Korea and the
US in April 2019. The main design decisions, each a response to an LTE limitation or a new
requirement, are worth listing before the details.
\begin{enumerate}
\item \textbf{One air interface from 410\,MHz to 71\,GHz.} Instead of one numerology, a family,
scaled by powers of two (Section~\ref{sec:ch21:numerology}).
\item \textbf{Ultra-lean design.} No always-on cell-specific reference signals; reference
signals are sent only when and where needed, which saves energy, reduces interference in
lightly loaded networks and leaves room for future features. The only always-on signal is the
synchronisation block, typically every 20\,ms.
\item \textbf{Beam-centric design.} Every signal and channel, including the broadcast ones, can
be beamformed, which is essential at millimetre waves and helps everywhere.
\item \textbf{Forward compatibility.} Resources can be reserved and left blank for future
features, and signals are self-contained in time and frequency where possible.
\item \textbf{Low latency by design.} Short slots at high numerologies, mini-slots, fast HARQ
feedback within the same slot when the device can manage it, and flexible TDD patterns.
\item \textbf{New channel coding.} LDPC for data and polar codes for control, replacing turbo
and convolutional codes (Section~\ref{sec:ch21:nrcoding}).
\end{enumerate}

\subsection{Flexible numerology, bandwidth parts and spectrum}
\label{sec:ch21:numerology}

NR's subcarrier spacing is $\Delta f=15\cdot2^{\mu}$\,kHz with $\mu=0,1,\ldots$, and a slot is
always 14 OFDM symbols (normal prefix), so slots last $1/2^{\mu}$\,ms: 1\,ms at 15\,kHz, 0.5\,ms at
30\,kHz, 0.125\,ms at 120\,kHz (TS~38.211). Chapter~\ref{ch:ofdm} derives the numerology table
and explains why the prefix fraction is constant and why symbols of different numerologies
nest (Section~\ref{sec:ch17:design}). The choice in practice follows the band:
\begin{itemize}
\item 15 or 30\,kHz in low bands (below about 2.5\,GHz), where cells are large and delay spreads
long, and where NR often shares a carrier with LTE;
\item 30\,kHz in the mid-band TDD bands around 3.5\,GHz, where most 5G capacity is, with
100\,MHz carriers of 273 RBs;
\item 120\,kHz in the millimetre-wave bands (FR2), with up to 400\,MHz per carrier, where phase
noise and Doppler demand wide spacing and cells are small.
\end{itemize}
The spectrum itself is divided into two \term{frequency ranges}: FR1 (410\,MHz to 7.125\,GHz)
and FR2 (24.25 to 52.6\,GHz, extended to 71\,GHz in Release~17). Figure~\ref{fig:ch21_spectrum_bands}
places the main cellular bands of every generation on one axis.

\mdcfig{ch21_spectrum_bands}{Major cellular bands by the generation that first used them, on a
logarithmic frequency axis (approximate edges, uplink and downlink shown together). Each
generation added spectrum above the last, but every band is eventually refarmed to the
newest technology; 5G uses everything from 600\,MHz to 40\,GHz. FR1 and FR2 are the 5G NR
frequency ranges of TS 38.101.}

A phone does not have to process the whole carrier. A \term{bandwidth part} (BWP) is a
contiguous set of RBs with its own numerology; a phone is configured with up to four per
direction and uses one at a time, switching by a DCI. A smartphone can sit on a 20\,MHz BWP
while browsing lightly and switch to the full 100\,MHz for a download, saving receiver power;
a low-cost device can operate only on a narrow part of a wide carrier.

\begin{inpractice}[Spectrum strategy: low, mid and high band]
Operators think about 5G spectrum in three layers. \textbf{Low band} (600--900\,MHz) gives
coverage and building penetration but only 5--20\,MHz per operator; NR here is often deployed
with \emph{dynamic spectrum sharing} (DSS), in which NR and LTE share a carrier and NR's data
channel is rate-matched around LTE's CRS. Low-band 5G is barely faster than LTE.
\textbf{Mid band} (2.5--4.2\,GHz, especially the ``C-band'' 3.3--4.2\,GHz, bands n77 and n78)
is the 5G sweet spot: 60--100\,MHz carriers per operator, TDD, massive MIMO, and cells that can
reuse existing macro sites with similar coverage to LTE at 1.8--2.1\,GHz once beamforming gain
is counted. In the US, the FCC's 2021 auction of 280\,MHz of C-band raised about \$81~billion,
and the 3.55--3.7\,GHz CBRS band introduced a three-tier licensing model in which a spectrum
access system database coordinates incumbents, priority licensees and general users.
\textbf{High band} (24--40\,GHz millimetre waves) offers 400--800\,MHz per operator and
multi-gigabit rates, but cells of a hundred metres or so, line-of-sight sensitivity and blockage
by bodies and foliage; it has been used mainly for stadiums, dense downtown areas and fixed
wireless access. Most 5G traffic in 2025 is carried in mid band.
\end{inpractice}

\subsection{Slot structure, TDD and mini-slots}

Each OFDM symbol in an NR slot can be downlink, uplink or \emph{flexible}, according to one of
a few dozen \emph{slot formats} (TS~38.213) combined with semi-static TDD patterns. Mid-band
networks typically use a fixed pattern coordinated among all operators in a country, because
unsynchronised TDD networks in adjacent blocks would jam each other's receivers. A common
choice at 30\,kHz is \texttt{DDDSU}: three downlink slots, a special slot (for example ten
downlink symbols, two guard symbols, two uplink symbols) and an uplink slot, repeating every
2.5\,ms (Figure~\ref{fig:ch21_nr_frame}c). That gives about three quarters of the time to the
downlink, matching traffic asymmetry, at the price of uplink latency: an uplink packet waits up
to 2.5\,ms for an uplink opportunity.

For latency-critical traffic, NR allows transmissions shorter than a slot. A
\term{mini-slot} (formally PDSCH/PUSCH ``mapping type~B'') of 2, 4 or 7 symbols (any length from
Release~16) can start at any symbol, and a URLLC transmission can \emph{pre-empt} an ongoing
eMBB transmission, with a pre-emption indication telling the eMBB user which symbols to
discard. Configured grants let a device transmit periodic uplink traffic without first
sending a scheduling request.

\mdcfig[0.95]{ch21_nr_frame}{NR frame elements (TS 38.211, TS 38.213). (a) The SS/PBCH block
(SSB): four OFDM symbols by 240 subcarriers, with the PSS and SSS on 127 subcarriers and the
PBCH (with its own DMRS) around them. (b) An SSB burst for case C (30\,kHz, 3--6\,GHz): up to
eight SSBs in the first half of a 5\,ms half-frame, each sent on a different beam, so a phone
anywhere in the sector hears at least one. (c) The DDDSU TDD pattern common in the 3.5\,GHz
band, and a two-symbol mini-slot that can start mid-slot to carry urgent traffic.}

\subsection{Initial access and beam management}

The NR \term{SS/PBCH block} (SSB) bundles the PSS, SSS and PBCH into four OFDM symbols and
240 subcarriers (20 RBs) (Figure~\ref{fig:ch21_nr_frame}a). The PSS and SSS (length-127
sequences, analysed in Section~\ref{sec:ch10:frame}) give 1008 physical cell identities;
the PBCH carries the master information block, which tells the phone where to find the
control resource set (CORESET\#0) that schedules the first system information block (SIB1).
Unlike LTE's PSS/SSS, which sit at the carrier centre, an SSB can be placed on a frequency
raster away from the centre, and a wide carrier may contain several.

The key change is that SSBs are \emph{beam-swept}. A cell transmits a burst of up to 4 SSBs
(below 3\,GHz), 8 (3--7\,GHz) or 64 (FR2) in a 5\,ms window, each in a different beam direction
(Figure~\ref{fig:ch21_nr_frame}b). The phone measures them, picks the best, and sends its random
access preamble in a PRACH occasion associated with that SSB index, which tells the base station
which beam to use for the reply. In connected mode, \term{beam management} refines the choice
with CSI reference signals swept in narrower beams and with the phone sweeping its own receive
beams (in FR2), reports the best beams, and recovers if the serving beam is blocked
(\emph{beam failure recovery}). Transmission configuration indicator (TCI) states tell the phone
which reference signal a given transmission is ``quasi-co-located'' with, that is, which beam
to expect. Chapter~\ref{ch:mimo} develops the array theory behind this.

\subsection{Physical channels}

NR's channels look familiar but are leaner and more flexible than LTE's.
\begin{itemize}
\item \textbf{PDCCH} is carried in configurable \term{control resource sets} (CORESETs) of
1--3 symbols and any set of RBs, not across the whole band, and phones search configured
\emph{search spaces} for their DCI. There is no PCFICH (the CORESET is configured) and no
PHICH: uplink retransmissions are scheduled explicitly, and HARQ is asynchronous in both
directions, with up to 16 processes and feedback timing indicated in the DCI.
\item \textbf{PDSCH/PUSCH} use front-loaded DMRS (in the third or fourth symbol of a slot, or the
first symbol of a mini-slot) so that the receiver can estimate the channel and start decoding
early, with optional additional DMRS symbols for high speed. Phase-tracking reference signals
(PT-RS) correct the common phase error caused by oscillator phase noise, a problem that grows
with carrier frequency (Chapter~\ref{ch:sdr}). Retransmission can be by \emph{code-block group},
so that one corrupted code block in a large transport block does not force retransmission of
all of them.
\item \textbf{The uplink waveform} is CP-OFDM by default, with DFT-spread OFDM available for
single-layer transmission when coverage matters (Section~\ref{sec:ch17:dfts}).
\item \textbf{Reference signals}: CSI-RS for channel measurement, beam management and fine time
and frequency tracking (the tracking reference signal, TRS), SRS for uplink sounding and TDD
reciprocity, and positioning reference signals.
\item \textbf{PRACH} uses long (839) or short (139) Zadoff--Chu sequences, the short ones for
the higher numerologies and small cells.
\end{itemize}
MCS tables go up to 256QAM (and 1024QAM on the FR1 downlink from Release~17), and there is a
low-spectral-efficiency table for URLLC, where a 10\% initial block error rate would be
unacceptable.

\subsection{Channel coding: LDPC and polar}
\label{sec:ch21:nrcoding}

NR replaced LTE's turbo code with \term{LDPC codes} for data and its convolutional code with
\term{polar codes} for control information (TS~38.212). Chapter~\ref{ch:moderncodes} develops
both; the choice itself was one of the most contested decisions in 3GPP's history.

For data, the requirements were throughputs of tens of gigabits per second, a wide range of
block lengths and code rates, and incremental redundancy. LDPC codes decode in parallel across
the whole codeword, whereas turbo decoding is iterative and partly sequential, so LDPC wins on
throughput per unit area and energy at high rates. NR uses two quasi-cyclic base graphs:
base graph~1 (for large blocks and high rates, up to 8448 information bits, lowest rate 1/3)
and base graph~2 (for small blocks and low rates, up to 3840 bits, rate down to 1/5), lifted
by circulant sizes up to 384. Both have a ``raptor-like'' structure that supports
incremental redundancy through a circular buffer and redundancy versions, as in LTE. The LDPC
decision was taken at the RAN1 meeting in Lisbon in October 2016.

For control, blocks are short (tens to a few hundred bits) and must be decoded reliably at low
SNR. After the next meeting, in Reno in November 2016, polar codes with CRC-aided successive
cancellation list decoding were adopted for downlink and uplink control and the broadcast
channel. Polar codes, invented by Erdal Ar\i{}kan in 2009, were the first codes provably
achieving capacity with low complexity; NR made them commercial within a decade. Very short
uplink control messages (up to 11 bits) use Reed--Muller or repetition codes.

\subsection{The peak rate formula}

TS~38.306 gives an approximate maximum data rate for a phone, used for capability signalling.
For a single carrier it is
\begin{equation}
R_{\max}\approx\nu_{\text{layers}}\,Q_m\,f\,R_{\max}^{\text{code}}\,
\frac{12\,N_{\text{PRB}}}{T_s^{\mu}}\,(1-\mathrm{OH}),
\qquad T_s^{\mu}=\frac{10^{-3}}{14\cdot2^{\mu}}\ \mathrm{s},
\label{eq:ch21:nrpeak}
\end{equation}
where $\nu$ is the number of MIMO layers, $Q_m$ the bits per modulation symbol, $f\le1$ a
scaling factor, $R_{\max}^{\text{code}}=948/1024$ the highest code rate, $N_{\text{PRB}}$ the
number of RBs, $T_s^{\mu}$ the average OFDM symbol duration including the prefix and OH the
overhead for control and reference signals: 0.14 for the FR1 downlink, 0.18 for the FR2 downlink,
0.08 and 0.10 for the uplink. For carrier aggregation the rates of the carriers are added.

\begin{worked}[NR peak rates in mid band and millimetre wave]
\emph{FR1, 100\,MHz at 30\,kHz.} $N_{\text{PRB}}=273$, $\mu=1$, so $T_s=35.7\,\mu$s and
$12\times273/T_s=91.7\times10^6$ resource elements per second. With four layers, 256QAM
($Q_m=8$), $f=1$:
\[
R=4\times8\times0.926\times91.7\times10^6\times(1-0.14)\approx2.34\ \text{Gb/s}.
\]
This counts every symbol as downlink. With the \texttt{DDDSU} pattern only about 74\% of
symbols are downlink, so the peak downlink rate of a TDD carrier is about 1.7\,Gb/s.

\emph{FR2, 400\,MHz at 120\,kHz.} $N_{\text{PRB}}=264$, $\mu=3$, $T_s=8.93\,\mu$s,
$12\times264/T_s=354.8\times10^6$ REs per second. With two layers and 64QAM, as typical for a
millimetre-wave phone:
\[
R=2\times6\times0.926\times354.8\times10^6\times(1-0.18)\approx3.23\ \text{Gb/s}.
\]
Aggregating two such carriers (800\,MHz) and adding mid-band carriers gives the 5--10\,Gb/s
peaks quoted for 2020s modems. Compare the ITU requirement of 20\,Gb/s: it is met by the
specification (with, say, eight layers and wide aggregation), not by any phone.
\end{worked}

\subsection{Massive MIMO in practice}

The biggest physical change at a 5G site is the antenna. A typical mid-band macro sector uses an
\term{active antenna unit} (AAU) with 64 transmit and receive chains (``64T64R'') driving an
array of about 192 dual-polarised elements, a panel roughly a metre tall and half a metre wide
with the radios and beamforming built in, radiating around 200\,W. Cheaper variants use 32 or
16 chains, or 8 for less loaded sites. In TDD, the base station estimates the downlink channel
from the phones' sounding reference signals by reciprocity; in FDD, or for phones whose
sounding is weak, it relies on CSI-RS and codebook feedback, either the coarse ``Type~I''
codebook or the high-resolution ``Type~II'' codebook that describes a few strongest paths with
their amplitudes and phases. The array delivers two gains (Chapter~\ref{ch:mimo}): about
15--20\,dB of beamforming gain, which lets 3.5\,GHz cells match the coverage of 1.8\,GHz LTE cells,
and multi-user MIMO, serving several phones on the same resource blocks with separate beams.
Published vendor and operator results have generally claimed capacity gains of a few times
over an 8-antenna baseline in busy cells, well short of the order-of-magnitude gains of
idealised analysis, because real users cluster, real channels are less ``rich'' than the
models, and channel information ages.

At millimetre waves, the arrays use \emph{hybrid} beamforming: a few RF chains each driving a
phased sub-array of tens of elements, at both ends. A phone contains several small antenna
modules on different edges, because a hand over one of them can block it.

\subsection{Non-standalone and standalone deployment}

To reach the market quickly, the first NR networks were \term{non-standalone} (NSA): NR
carriers added as a secondary node to an LTE connection, using \term{E-UTRA--NR dual
connectivity} (EN-DC) and the existing EPC (``option~3'' in 3GPP's deployment-option numbering,
with variants 3a and 3x for how user traffic is split). The LTE eNB anchors the control
signalling and mobility, and the NR gNB adds capacity. NSA needed no new core and reused LTE
coverage, but it delivered only eMBB and inherited LTE's control-plane latency. In
\term{standalone} (SA, ``option~2'') operation, NR connects to the new 5G core with no LTE
involvement, enabling the 5G core's features: network slicing, the RRC\_INACTIVE state, lower
latency and new voice options (VoNR). Figure~\ref{fig:ch21_nsa_sa} compares them.

\begin{figure}[htbp]\centering
\begin{tikzpicture}[font=\sffamily\scriptsize,
  box/.style={draw=navy,fill=navy!6,thick,rounded corners=1pt,minimum height=0.7cm,minimum width=1.3cm,align=center},
  nr/.style={draw=accent,fill=accent!8,thick,rounded corners=1pt,minimum height=0.7cm,minimum width=1.3cm,align=center},
  core/.style={draw=practc,fill=practc!8,thick,rounded corners=2pt,minimum height=0.7cm,minimum width=2.6cm,align=center},
  cp/.style={thick,black!70,dashed}, up/.style={thick,practc}]
\node[font=\sffamily\bfseries\scriptsize,text=navy] at (2.2,3.0) {Option 3x: non-standalone (EN-DC)};
\node[core] (epc) at (2.2,2.2) {EPC (MME, S-GW, P-GW)};
\node[box] (enb) at (0.9,0.8) {eNB\\(LTE anchor)};
\node[nr] (gnb) at (3.5,0.8) {en-gNB\\(NR secondary)};
\node[box] (ue) at (2.2,-0.6) {UE (LTE + NR)};
\draw[cp] (enb) -- node[left,font=\sffamily\tiny]{S1-MME} (epc);
\draw[up] (gnb) -- node[right,font=\sffamily\tiny]{S1-U} (epc);
\draw[thick,navy] (enb) -- node[above,font=\sffamily\tiny]{X2} (gnb);
\draw[thick,navy,dashed] (ue) -- (enb); \draw[thick,accent,dashed] (ue) -- (gnb);
\node[font=\sffamily\bfseries\scriptsize,text=navy] at (8.7,3.0) {Option 2: standalone};
\node[core] (5gc) at (8.7,2.2) {5GC (AMF, SMF, UPF)};
\node[nr] (g1) at (7.5,0.8) {gNB};
\node[nr] (g2) at (9.9,0.8) {gNB};
\node[box] (ue2) at (8.7,-0.6) {UE (NR only)};
\draw[cp] (g1) -- node[left,font=\sffamily\tiny]{N2} (5gc); \draw[up] (g2) -- node[right,font=\sffamily\tiny]{N3} (5gc);
\draw[thick,navy] (g1) -- node[above,font=\sffamily\tiny]{Xn} (g2);
\draw[thick,accent,dashed] (ue2) -- (g1);
\draw[navy!40,thick] (5.6,-0.9) -- (5.6,3.2);
\end{tikzpicture}
\caption{Non-standalone (left) and standalone (right) 5G. In option 3x the LTE eNB carries the
control signalling and the NR node adds user-plane capacity through dual connectivity, all
attached to the 4G core. In option 2, NR attaches to the 5G core and works alone. Dashed black
lines: control; green: user plane; dashed coloured lines: radio links.}
\label{fig:ch21_nsa_sa}
\end{figure}

\begin{pitfall}[What the ``5G'' icon means]
The 5G indicator on a phone says very little. In NSA operation it may appear when the phone is
attached to an LTE anchor cell that \emph{could} add an NR carrier, even if no NR carrier is
active. A low-band NR carrier shared with LTE by DSS can be slower than the LTE it displaced.
And for a few years one US operator labelled its LTE-Advanced service ``5G~E''. When you
measure or compare 5G performance, record the band, bandwidth, numerology, NSA or SA mode and
the number of carriers and layers actually in use (all visible in engineering modes and drive-test
tools), not the icon.
\end{pitfall}

\subsection{The 5G core: a service-based architecture}
\label{sec:ch21:5gc}

The \term{5G core} (5GC, 3GPP TS~23.501 and 23.502) was redesigned as software. Instead of
boxes connected by point-to-point protocols, its control plane is a set of
\term{network functions} (NFs) that offer services to each other over a common
\term{service-based interface} using HTTP/2 and JSON, as cloud applications do
(Figure~\ref{fig:ch21_sba}). The main functions are:
\begin{itemize}
\item the \term{access and mobility management function} (AMF): registration, connection and
mobility management, the endpoint of NAS signalling (part of LTE's MME);
\item the \term{session management function} (SMF): PDU sessions, IP address allocation and
control of the user-plane functions (the rest of the MME and the control part of the gateways);
\item the \term{user plane function} (UPF): packet forwarding, QoS enforcement, usage reporting,
the anchor for mobility and the connection to data networks, controlled by the SMF over N4
(the S-GW and P-GW user planes, separated as in LTE's later ``CUPS'' architecture);
\item the authentication server function (AUSF), the unified data management (UDM) and data
repository (UDR), the policy control function (PCF), the network repository function (NRF,
where NFs register and discover each other), the network slice selection function (NSSF),
the network exposure function (NEF, which exposes capabilities to outside applications) and
application functions (AF).
\end{itemize}
Because UPFs are separate and lightweight, they can be placed anywhere: centrally for ordinary
internet traffic, or at the edge of the network for low latency (\term{multi-access edge
computing}, MEC), where an application server sits next to a local UPF a few kilometres from the
base station instead of hundreds of kilometres away.

\begin{figure}[htbp]\centering
\begin{tikzpicture}[font=\sffamily\scriptsize,
  nf/.style={draw=navy,fill=navy!6,thick,rounded corners=2pt,minimum height=0.62cm,minimum width=1.05cm,align=center},
  up/.style={draw=practc,fill=practc!8,thick,rounded corners=2pt,minimum height=0.7cm,minimum width=1.2cm,align=center},
  ra/.style={draw=accent,fill=accent!8,thick,rounded corners=2pt,minimum height=0.7cm,minimum width=1.2cm,align=center}]
\draw[thick,navy] (0.0,2.2) -- (10.8,2.2);
\node[font=\sffamily\tiny,text=navy,anchor=west] at (10.85,2.2) {SBI};
\foreach \x/\nm in {0.6/NSSF,1.95/NEF,3.3/NRF,4.65/PCF,6.0/UDM,7.35/AF,8.7/AUSF}{
  \node[nf] (\nm) at (\x,3.0) {\nm}; \draw[thick,navy] (\nm.south) -- (\x,2.2);}
\node[nf] (AMF) at (2.6,1.35) {AMF}; \draw[thick,navy] (AMF.north) -- (2.6,2.2);
\node[nf] (SMF) at (6.4,1.35) {SMF}; \draw[thick,navy] (SMF.north) -- (6.4,2.2);
\node[nf] (UDR) at (10.0,3.0) {UDR}; \draw[thick,navy] (UDR.south) -- (10.0,2.2);
\node[ra] (ue) at (0.2,-0.6) {UE};
\node[ra] (gnb) at (2.6,-0.6) {(R)AN\\gNB};
\node[up] (upf) at (6.4,-0.6) {UPF};
\node[draw=gray,thick,ellipse,minimum width=1.5cm] (dn) at (9.4,-0.6) {DN};
\draw[thick,black!60,dashed] (ue.north) to[bend left=20] node[left,font=\sffamily\tiny]{N1} (AMF.west);
\draw[thick,black!60,dashed] (gnb) -- node[left,font=\sffamily\tiny]{N2} (AMF);
\draw[thick,black!60,dashed] (SMF) -- node[right,font=\sffamily\tiny]{N4} (upf);
\draw[thick,practc] (ue) -- (gnb); \draw[thick,practc] (gnb) -- node[above,font=\sffamily\tiny]{N3} (upf);
\draw[thick,practc] (upf) -- node[above,font=\sffamily\tiny]{N6} (dn);
\draw[thick,practc] (upf.south) .. controls +(0.4,-0.6) and +(-0.4,-0.6) .. node[below,font=\sffamily\tiny]{N9 (to another UPF)} ([xshift=6mm]upf.south);
\end{tikzpicture}
\caption{The 5G core's service-based architecture (TS 23.501). Control-plane network
functions offer services to each other over a common service-based interface (SBI, HTTP/2);
only the radio-facing and user-plane links (N1--N4, N6, N9) remain point-to-point. The UPF can
be instantiated centrally or at the network edge.}
\label{fig:ch21_sba}
\end{figure}

\subsubsection{Sessions, QoS flows and slices}

A device obtains connectivity through a \term{PDU session} to a data network, identified by a
data network name (the successor of the APN). Within a session, packets are classified into
\term{QoS flows}, each with a 5G QoS identifier (5QI) that sets its resource type (guaranteed
bit rate, delay-critical or best effort), priority, delay budget and error-rate target. A new
radio protocol layer above PDCP, SDAP, maps QoS flows to radio bearers.

\textbf{Network slicing}\index{network slicing} lets an operator run several logically separate networks on the same
infrastructure, each with its own set of NF instances, policies and possibly reserved radio
resources. A slice is identified by an S-NSSAI: an 8-bit slice/service type (1 for eMBB,
2 for URLLC, 3 for massive IoT, 4 for V2X, among others) and an optional 24-bit slice
differentiator. A factory, an emergency service or a mobile virtual network operator can have a
slice with guaranteed resources and isolated management. Slicing was one of 5G's most
heavily promoted features; its commercial use grew slowly, because the radio part of a slice
(guaranteeing resources in a shared cell) and the business models proved harder than the core
part.

\subsubsection{Security}

5G kept LTE's mutual authentication and added two important fixes. The permanent identity
(SUPI, the successor of the IMSI) is never sent in the clear: the phone encrypts it with the
home network's public key (an elliptic-curve integrated encryption scheme) into a
\emph{subscription concealed identifier} (SUCI), which defeats passive IMSI catchers. And the
home network now confirms authentication itself (5G-AKA or EAP-AKA$'$), rather than trusting the
visited network, closing a roaming fraud gap. 5G also added integrity protection for user-plane data, which LTE
lacked (its use is set by network policy), in response to attacks demonstrated on LTE that
modified unprotected user packets in flight.

\subsection{The disaggregated RAN: CU, DU, RU and O-RAN}

The 5G base station, the \term{gNB}, can be split into parts (Figure~\ref{fig:ch21_ran_split}).
The \term{central unit} (CU) runs the non-real-time upper layers (RRC, SDAP, PDCP) and can be
centralised and virtualised on ordinary servers; it can itself be split into control-plane
and user-plane parts. The \term{distributed unit} (DU) runs the real-time layers (RLC, MAC and
the upper physical layer: coding, modulation, layer mapping), close to the antennas. The
\term{radio unit} (RU) does the lower physical layer (FFT, cyclic prefix, digital beamforming,
digital-to-analog conversion) and the RF. The CU--DU interface F1 is standardised by 3GPP (the
``option~2'' split between PDCP and RLC); the DU--RU \emph{fronthaul} carries frequency-domain IQ
samples over Ethernet (eCPRI), using the ``option~7.2x'' split specified by the
\term{O-RAN Alliance}. Splitting the PHY that way reduces the fronthaul rate by sending
frequency-domain samples only for the occupied subcarriers, while keeping the heavy, low-latency
processing in the DU.

\begin{figure}[htbp]\centering
\begin{tikzpicture}[font=\sffamily\scriptsize,
  box/.style={draw=navy,fill=navy!6,thick,rounded corners=2pt,minimum width=2.4cm,align=center},
  ric/.style={draw=accent,fill=accent!6,thick,rounded corners=2pt,minimum width=2.4cm,align=center},
  arr/.style={thick,navy,{Stealth[length=1.6mm]}-{Stealth[length=1.6mm]}}]
\node[box,minimum height=1.6cm] (ru) at (0,0) {\textbf{RU}\\[2pt]low PHY: FFT/iFFT,\\CP, digital\\beamforming, RF};
\node[box,minimum height=1.6cm] (du) at (4.0,0) {\textbf{DU}\\[2pt]high PHY: coding,\\modulation, MIMO\\RLC, MAC (scheduler)};
\node[box,minimum height=1.6cm] (cu) at (8.0,0) {\textbf{CU} (CU-CP, CU-UP)\\[2pt]RRC, SDAP, PDCP\\virtualised,\\centralised};
\node[draw=practc,fill=practc!8,thick,rounded corners=2pt,minimum width=1.7cm,minimum height=0.8cm] (core) at (11.3,0) {5G core};
\draw[arr] (ru) -- node[above,font=\sffamily\tiny]{open fronthaul} node[below,font=\sffamily\tiny]{split 7.2x, eCPRI} (du);
\draw[arr] (du) -- node[above,font=\sffamily\tiny]{F1} node[below,font=\sffamily\tiny]{midhaul} (cu);
\draw[arr] (cu) -- node[above,font=\sffamily\tiny]{N2/N3} node[below,font=\sffamily\tiny]{backhaul} (core);
\node[ric,minimum height=0.7cm] (nrt) at (6.0,1.8) {near-RT RIC (10\,ms--1\,s)\\xApps};
\node[ric,minimum height=0.7cm] (srt) at (10.4,1.8) {non-RT RIC ($>$1\,s)\\rApps, in SMO};
\draw[thick,accent,dashed] (nrt) -- node[left,font=\sffamily\tiny]{E2} (du.north);
\draw[thick,accent,dashed] (nrt) -- (cu.north);
\draw[thick,accent,dashed] (srt) -- node[above,font=\sffamily\tiny]{A1} (nrt);
\node[font=\sffamily\tiny,text=black!60] at (0,-1.15) {at the antenna};
\node[font=\sffamily\tiny,text=black!60] at (4.0,-1.15) {at the site or a hub ($<$\,20\,km)};
\node[font=\sffamily\tiny,text=black!60] at (8.0,-1.15) {in an edge data centre};
\end{tikzpicture}
\caption{The disaggregated 5G RAN. 3GPP standardises the CU--DU (F1) interface; the O-RAN
Alliance specifies the open fronthaul between DU and RU and the RAN intelligent controllers
(RICs), which host third-party applications that tune the RAN on near-real-time and
non-real-time scales.}
\label{fig:ch21_ran_split}
\end{figure}

The O-RAN Alliance, formed by operators in 2018, aims to make these interfaces open so that
RUs, DUs and CUs from different vendors interoperate, and to add \term{RAN intelligent
controllers} (RICs) that run third-party applications (xApps and rApps) for functions such as
traffic steering, energy saving and interference management. Open RAN is the latest round of a
recurring argument in cellular---GSM's open Abis interface was supposed to do the same in
1991---and its success depends as much on integration cost and operator commitment as on the
specifications.

\subsection{Releases 15 to 18: 5G grows beyond broadband}

Release~15 delivered eMBB. The following releases turned NR into a platform for other
industries and filled the gaps (the dates are the 3GPP functional freezes).
\begin{itemize}
\item \textbf{Release~16} (2020): URLLC enhancements (multiple HARQ feedbacks per slot,
uplink pre-emption, PDCP duplication over two legs, multi-TRP transmission), \emph{industrial
IoT} with time-sensitive networking integration, positioning with downlink and uplink time
difference, multi-cell round-trip time and angle methods (targets of a few metres),
NR in unlicensed spectrum (NR-U), integrated access and backhaul (IAB), NR sidelink for V2X,
two-step random access and device power saving.
\item \textbf{Release~17} (2022): \term{RedCap} (reduced-capability) devices for wearables,
industrial sensors and video surveillance, with 20\,MHz bandwidth in FR1 and one or two
receive antennas; \term{non-terrestrial networks} (NTN), adapting NR and NB-IoT to satellite
links with long delays and large Doppler shifts (Chapter~\ref{ch:satellite},
Section~\ref{sec:ch23:ntn}); operation up to 71\,GHz; positioning with sub-metre accuracy targets
for industrial uses; multicast and broadcast; coverage enhancements; and further URLLC and
sidelink work.
\item \textbf{Release~18} (2024), the first \term{5G-Advanced} release: studies of artificial
intelligence and machine learning for the air interface (CSI compression, beam management,
positioning), network energy saving, extended reality, enhanced RedCap with 5\,MHz bandwidth,
network-controlled repeaters, mobile IAB, low-power wake-up signals, and NTN enhancements.
Releases~19 and~20 continue 5G-Advanced and begin the studies for 6G, the subject of the final
chapter.
\end{itemize}

\subsection{Latency: a budget, not a number}

``5G latency is 1\,ms'' is one of the most misquoted claims in telecommunications. The ITU
requirement is for the \emph{user-plane air-interface} latency of a small packet, one way,
in an unloaded system. Figure~\ref{fig:ch21_latency} puts it in context.

\mdcfig{ch21_latency}{(a) Typical round-trip times measured by users of each technology to a
nearby server (approximate ranges; real values depend on load, core-network routing and server
location). (b) An illustrative one-way downlink air-interface latency budget in the style of the
IMT-Advanced and IMT-2020 self-evaluations: waiting for the next transmission opportunity
(alignment), the transmission time, base-station and device processing, and the expected cost
of a HARQ retransmission with 10\% probability.}

\begin{worked}[A latency budget]
\emph{LTE FDD.} A packet arriving at the eNB waits on average half a subframe for the next one
(0.5\,ms), is processed and scheduled (about 1\,ms), transmitted (1\,ms TTI) and decoded by the phone
(about 1.5\,ms). With a 10\% chance of one retransmission costing 8\,ms, the expected
budget is $0.5+1+1+1.5+0.1\times8=4.8$\,ms, just inside LTE's 5\,ms target. The uplink is worse:
a phone must first send a scheduling request and wait for a grant, adding several
milliseconds, unless it has a pre-configured grant.

\emph{NR, 30\,kHz.} The slot is 0.5\,ms. Scaling the processing times with the slot (devices of
NR ``processing capability~1'' need roughly 10 symbols from reception to HARQ feedback at
30\,kHz) gives alignment 0.25, transmission 0.5, processing about 1\,ms in total and HARQ
$0.1\times2=0.2$\,ms: about 2\,ms.

\emph{NR URLLC, 120\,kHz with a two-symbol mini-slot.} The transmission is
$2\times8.9\,\mu\mathrm{s}\approx18\,\mu$s, alignment half that, and with fast (``capability 2'')
processing the total is a few tenths of a millisecond, comfortably inside 1\,ms.

\emph{End to end.} Add the transport network (a few milliseconds per hundred kilometres of fibre
plus switching), the core (a fraction of a millisecond in a well-placed UPF) and the server.
A user in a city reaching a server in a distant data centre sees 20--40\,ms on LTE and 10--25\,ms on
5G; a device on a factory network with a local UPF and an edge server can see a few milliseconds.
The radio is now rarely the largest term.
\end{worked}

\begin{inpractice}[5G in the field versus the brochure]
The first five years of 5G showed a familiar gap between specification and experience.
Mid-band 5G with 100\,MHz carriers and massive MIMO delivered large real gains: median
downlink speeds measured by crowd-sourced testing services were typically several times those
of LTE in the same markets, often in the low hundreds of megabits per second, with peaks above
1\,Gb/s. Low-band and DSS deployments gave little improvement over LTE. Millimetre-wave
networks delivered gigabit speeds in line of sight of a node and nothing a few buildings away.
Uplink speeds and latency improved less than downlink speeds, because TDD patterns favour the
downlink and the core and transport networks dominate latency. The eMBB story was largely
delivered; URLLC and massive IoT found their first markets in private industrial networks and
in NB-IoT and LTE-M rather than in the public network. The lesson for an engineer is to read the
peak-rate formula, the TDD pattern and the deployment band before believing any number.
\end{inpractice}

%======================================================================
\section{The generations side by side}
\label{sec:ch21:compare}
%======================================================================

Table~\ref{tab:ch21:compare} lines up the five generations, and
Figure~\ref{fig:ch21_peak_rates} plots their peak rates against time. Over three decades the
peak rate available to a phone grew from 9.6\,kb/s to several gigabits per second, a factor of
about $10^6$: roughly a factor of ten every five years. It is instructive to decompose that
factor. Bandwidth grew from 200\,kHz (GSM) to several hundred megahertz of aggregated carriers,
a factor of about 1000--4000. Peak spectral efficiency grew from about 0.05\,b/s/Hz per user
(a GSM data slot) to about 30\,b/s/Hz (NR with many layers and 256QAM), a factor of several
hundred, of which MIMO layers contribute a factor of four to eight, denser modulation a factor
of eight, and better coding and less overhead the rest. Bandwidth, in other words, did most of
the work, which is why spectrum is the most expensive input to a mobile network.

\mdcfig{ch21_peak_rates}{Peak downlink rates of each generation's standards and device
categories against the year they became commercially available (approximate). The dotted
line is a least-squares exponential fit: about a factor of ten every five years. Peak rates
are what a single user could get alone in a cell with an excellent channel; typical rates in
loaded networks are roughly a tenth of these.}

\begin{table}[htbp]\centering\footnotesize
\caption{The cellular generations compared. Peak rates are for the specification or for
commercial devices of the period, as noted; ``typical latency'' is a representative round-trip
time seen by users (approximate).}
\label{tab:ch21:compare}
\setlength{\tabcolsep}{3pt}
\begin{tabularx}{\linewidth}{@{}>{\raggedright\arraybackslash}p{2.0cm} *{5}{>{\raggedright\arraybackslash}X}@{}}
\toprule
 & \textbf{1G} & \textbf{2G / 2.5G} & \textbf{3G / 3.5G} & \textbf{4G} & \textbf{5G}\\
\midrule
Main systems & AMPS, NMT, TACS & GSM, IS-95, IS-136, PDC; GPRS, EDGE & WCDMA, CDMA2000, TD-SCDMA; HSPA, EV-DO & LTE, LTE-Advanced (Pro) & NR (Releases 15--18)\\
First service & 1979--83 & 1991--95 & 2001--03 & 2009--10 & 2019\\
Carrier bandwidth & 25--30\,kHz & 200\,kHz (GSM), 1.25\,MHz (IS-95) & 5\,MHz (WCDMA), 1.25\,MHz (CDMA2000) & 1.4--20\,MHz, up to 32 carriers aggregated & 5--100\,MHz (FR1), 50--400\,MHz (FR2)\\
Multiple access & FDMA & TDMA, CDMA & CDMA (+ TDM scheduling in HSPA/EV-DO) & OFDMA down, SC-FDMA up & OFDMA; CP-OFDM or DFT-s-OFDM up\\
Duplexing & FDD & FDD & FDD (TDD) & FDD and TDD & mostly TDD above 2.5\,GHz\\
Modulation & analog FM & GMSK, $\pi/4$-DQPSK, 8-PSK & QPSK to 64QAM & QPSK to 256QAM & up to 256QAM (1024QAM)\\
Channel coding & BCH (signalling) & convolutional & convolutional, turbo & turbo, TBCC & LDPC, polar\\
MIMO & -- & -- & $2\times2$ (HSPA+) & up to 8 layers DL & up to 8 layers DL; massive MIMO arrays, beam management\\
Peak rate (DL) & -- (analog voice) & 9.6\,kb/s (CSD) to 237\,kb/s (EDGE) & 384\,kb/s (R99) to 42--168\,Mb/s (HSPA+) & 150\,Mb/s (Rel-8 phone) to $\sim$1\,Gb/s (LTE-A) & 2--10\,Gb/s (devices); 20\,Gb/s (requirement)\\
Peak spectral efficiency & -- & $\sim$1.4\,b/s/Hz gross per carrier & $\sim$8\,b/s/Hz (HSPA+ MIMO) & 15\,b/s/Hz (IMT-Advanced) & 30\,b/s/Hz (IMT-2020)\\
Typical latency (RTT) & -- & 300--1000\,ms (GPRS/EDGE) & 100--200\,ms (R99), 30--100\,ms (HSPA) & 20--50\,ms & 8--25\,ms (eMBB)\\
Core network & analog switch & circuit (MSC) + GPRS packet core & circuit + packet (Iu-CS, Iu-PS) & all-IP EPC; voice over IMS & service-based 5GC, slicing, edge UPF\\
Defining idea & cellular reuse, handoff & digital voice, SIM, roaming, SMS & wideband CDMA, packet data & OFDMA, flat IP, scheduler-centric & flexible numerology, beams, cloud core\\
\bottomrule
\end{tabularx}
\end{table}

%======================================================================
\section{Lessons from forty years}
\label{sec:ch21:lessons}
%======================================================================

\subsection{What each generation got right and wrong}

\paragraph{1G} proved that cellular reuse and handoff worked and that people would pay for
mobility; it got the business right and the security disastrously wrong, and its national
fragmentation (outside North America and the Nordic countries) left travellers stranded.

\paragraph{2G (GSM)} got the ecosystem right: a single specification across many countries,
open interfaces, the SIM card, international roaming and a guaranteed market through the MoU.
Its radio design was conservative and robust---constant-envelope GMSK, an equaliser, slow
frequency hopping---and its features (SMS, prepaid, roaming) outlived several successors. Its
security relied on secret algorithms and one-way authentication, and its circuit-switched
architecture made data a retrofit. IS-95 got interference averaging, power control and soft
handoff right a decade before anyone else, and its main error was commercial: it never
achieved GSM's scale.

\paragraph{3G} got the technology roughly right and the timing and economics wrong. WCDMA was
an elegant, capable radio, but its architecture put the radio intelligence in the RNC, its
first releases were tuned for circuit-like services, and the operators had paid for spectrum
as though data revenue would arrive in 2002. It arrived in 2008, delivered by HSPA and the
smartphone. The lasting contributions were turbo codes, fast power control, the 3GPP process
itself and mutual authentication.

\paragraph{4G (LTE)} got almost everything right for its decade: OFDMA, a flat architecture,
a scheduler-centric design, one global standard, and enough headroom (carrier aggregation,
MIMO) to grow from 100\,Mb/s to a gigabit over ten years. Its regrets were always-on
cell-specific reference signals (an energy and interference cost that grows with the number of
cells, and an obstacle to beamforming), a single numerology, and a late, complicated path to
voice.

\paragraph{5G} fixed LTE's regrets and delivered eMBB well in mid band. It overpromised in
URLLC, network slicing and millimetre-wave coverage for consumers, at least on the timescale
marketed, and its business case for operators has been weaker than its technology. The
flexible, lean, beam-based air interface and the cloud-native core are likely to be its
longest-lived contributions, inherited by 6G.

\begin{keyidea}[Standards are bets about the future]
Every generation's designers had to guess, five to ten years ahead, what users would want and
what silicon could do. The ones that aged well---GSM, LTE---left headroom and kept options open:
GSM's signalling channels made SMS possible; LTE's flexible scheduler and carrier aggregation
absorbed a hundredfold growth. The ones that aged badly were tuned too precisely to a predicted
service (3G's circuit-switched video calls) or relied on assumptions that did not hold (secret
algorithms, always-on reference signals). When you design a system that must live for twenty
years, design for flexibility in the parts you are least sure of, and simplicity in the parts
that must never fail.
\end{keyidea}

\subsection{The economics}

Engineering decisions in cellular are inseparable from money. A few facts frame them.
\begin{itemize}
\item \textbf{Traffic grows, revenue does not.} Mobile data traffic has grown by tens of percent
a year for fifteen years, while operators' revenues in most mature markets have been flat in
real terms. The cost of carrying a bit must therefore fall at roughly the rate traffic grows,
and every generation must deliver a large reduction in cost per bit, not just higher rates.
Spectral efficiency, wider carriers, cheaper sites and lower energy are all ways to that end.
\item \textbf{Spectrum is the most expensive input.} Auctions have raised hundreds of billions
of dollars worldwide since 2000. Because bandwidth is the largest contributor to capacity, an
operator's spectrum holdings largely determine its network quality, and refarming old bands to
new technologies is a constant activity.
\item \textbf{Energy is the largest operating cost of the RAN.} A macro site draws several
kilowatts, much of it while carrying little traffic. NR's lean design, sleep modes and
Release~18's network energy saving features respond to an operating cost, not a technical
curiosity.
\item \textbf{Sites are expensive and slow.} Acquiring, building and powering a new site takes
months to years. This is why mid-band 5G was designed to reuse existing macro sites, why small
cells grew more slowly than predicted, and why tower companies, which own and lease the sites,
became some of the most valuable businesses in telecommunications.
\item \textbf{Scale wins.} The cost of a standard's chips and handsets falls with volume, so a
standard with a global market (GSM, LTE, NR) beats a technically comparable regional one
(D-AMPS, PDC, WiMAX, TD-SCDMA after its domestic phase). Standards-essential patents, licensed
on ``fair, reasonable and non-discriminatory'' terms, fund the research that produces each
generation and are a permanent source of litigation.
\end{itemize}

\begin{inpractice}[The ten-year rhythm and what comes next]
If the pattern holds, 6G networks will open around 2030. 3GPP began Release~20 studies in
2025, and the ITU's IMT-2030 framework (Recommendation ITU-R M.2160, 2023) adds integrated
sensing and communication, native AI and ubiquitous coverage (including non-terrestrial
networks) to the usual extrapolation of rates and latency. The candidate technologies---new
spectrum around 7--15\,GHz and above 100\,GHz, extremely large antenna arrays, new waveforms,
AI-native air interfaces---are the subject of Chapter~\ref{ch:sixg}. The lesson of this chapter is
that whichever technologies are chosen, the system will be judged on the same unglamorous
criteria as its predecessors: cost per bit, energy per bit, coverage, reliability, security and
the willingness of a whole industry to build it together.
\end{inpractice}

%======================================================================
\section*{Summary}
%======================================================================
\begin{itemize}
\item Cellular systems reuse frequencies in cells, hand calls over between them, and grow
capacity by shrinking cells. Every generation since 1G inherits this structure; what changes is
how users share a cell's spectrum and how interference between cells is managed.
\item 1G (AMPS, NMT, TACS) used analog FM in 25--30\,kHz channels with reuse of about 7,
and was undone by capacity limits, eavesdropping and cloning.
\item GSM introduced digital TDMA (eight slots on 200\,kHz carriers, 270.833\,kb/s GMSK), a
burst with a midamble for an MLSE equaliser, unequal error protection of speech, slow frequency
hopping, DTX, power control, the SIM and challenge--response authentication, open network
interfaces and international roaming. Its ciphers and one-way authentication were broken.
GPRS and EDGE added packet data, 8-PSK, link adaptation and incremental redundancy.
\item IS-95 and then WCDMA built capacity on interference averaging: universal reuse,
fast closed-loop power control (800 and 1500\,Hz), soft handover and variable-rate vocoders.
HSPA moved scheduling, HARQ and link adaptation into the base station with a 2\,ms TTI and
shared channels, the template for everything after.
\item LTE adopted OFDMA downlink and SC-FDMA uplink, a flat all-IP architecture with one radio
node, 1\,ms subframes and 180\,kHz resource blocks, turbo coding with circular-buffer
incremental redundancy, 8\,ms HARQ, CQI-driven link adaptation, MIMO transmission modes,
carrier aggregation and voice over IMS.
\item NR scales the numerology ($15\cdot2^{\mu}$\,kHz), removes always-on reference signals,
beamforms everything including synchronisation (SSB bursts), uses LDPC and polar codes, supports
mini-slots and flexible TDD, and runs on a service-based, cloud-native core with network
slicing and edge computing; the RAN can be split into CU, DU and RU with open interfaces.
\item Peak rates have grown about tenfold every five years, mostly through bandwidth; typical
rates are about a tenth of the peaks, and latency is a budget in which the radio is now a small
term.
\end{itemize}

\begin{labbox}[Cells, air interfaces and the LTE/NR physical layer]
\lab{moderncomms-labs/labs/lab11\_air\_interfaces.py} accompanies this chapter and
Chapters~\ref{ch:multipleaccess} and~\ref{ch:wifi}. It computes co-channel SIR against cluster
size (the AMPS and GSM reuse trade-off of Sections~\ref{sec:ch21:1g} and~\ref{sec:ch21:gsm}),
sizes cells with Erlang~B, tabulates the NR numerologies and draws a simplified NR slot with its
control region, DM-RS and PT-RS overhead, compares round-robin, max-rate and proportional-fair
OFDMA schedulers, measures the PAPR advantage of DFT-s-OFDM that justified LTE's SC-FDMA uplink,
builds a link-adaptation staircase from the NR MCS table with HARQ, and simulates Wi-Fi
contention. A planned follow-on lab, \lab{moderncomms-labs/labs/lab27\_lte\_nr\_phy.py}, will
build an LTE or NR downlink end to end: a resource grid with synchronisation signals and
reference signals as in Figures~\ref{fig:ch21_lte_grid} and~\ref{fig:ch21_nr_frame}, PSS/SSS
cell search on a simulated multi-cell signal, and a PDSCH chain with CRC, LDPC coding,
circular-buffer rate matching, QAM, OFDM and HARQ combining. The figure script
\lab{book/figscripts/ch21\_figs.py} reproduces every data figure in this chapter, including the
power-control and soft-handover simulations.
\end{labbox}

\problems
\begin{enumerate}[label=\thechapter.\arabic*]
\item An AMPS operator has 395 voice channels. Compare the number of erlangs it can carry at 2\%
blocking with (a) 7-cell omnidirectional reuse and (b) 7-cell reuse with three sectors per cell.
Explain why sectorisation improves the SIR but reduces trunking efficiency, and estimate the
net effect on subscribers per cell.
\item Starting from the 13\,MHz GSM master clock, derive the bit rate, bit period, timeslot
duration, TDMA frame duration and 26-multiframe duration. Why does the 26-multiframe last exactly
120\,ms?
\item Verify the GSM full-rate channel-coding arithmetic: show how 260 speech bits become 456
coded bits, and compute the code rate applied to the class 1 bits. If the class 2 bits had
also been coded at rate 1/2, what gross rate would a traffic channel need?
\item A GSM base station measures that a mobile's access burst arrives 74\,$\mu$s late. What
timing advance (in bit periods) should it command, and how far away is the mobile? What is the
maximum cell radius with the standard access burst guard period?
\item In a GSM network with slow frequency hopping over $N$ carriers, assume each burst
independently sees a deep fade with probability $p=0.1$ and the decoder fails if more than two
of the eight bursts of a speech frame are lost. Compute the frame erasure rate with hopping
(independent bursts) and without hopping for a stationary user whose eight bursts fade together.
\item For an EDGE user on four timeslots, compute the throughput with MCS-5, MCS-7 and MCS-9.
If the BLERs at the user's C/I are 1\%, 15\% and 60\% respectively and failed blocks are
retransmitted until they succeed (ignore incremental redundancy), which MCS gives the highest
goodput?
\item A WCDMA uplink serves voice users at 12.2\,kb/s with a required $\ebno$ of 5\,dB, a voice
activity factor of 0.6 and an other-cell to own-cell interference ratio of 0.65. Using the
classic single-cell CDMA capacity formula with these corrections, estimate the number of users
per cell for a noise rise of 3\,dB and of 6\,dB. Relate the difference to cell breathing.
\item Explain why HSDPA has no fast power control on its shared channel and no soft handover,
while WCDMA Release~99 depends on both. What replaces them?
\item Compute the number of resource elements per subframe available for the PDSCH in a
10\,MHz LTE carrier with two CRS ports and a two-symbol control region, in a subframe without
synchronisation or broadcast signals. What peak rate does that give with $2\times2$ MIMO and
64QAM at code rate 0.9?
\item Using~\eqref{eq:ch21:nrpeak}, compute the maximum NR downlink rate for (a) 20\,MHz at
15\,kHz (106 RBs), 4 layers, 256QAM, FDD; (b) 100\,MHz at 30\,kHz, 4 layers, 256QAM, with the
\texttt{DDDSU} pattern of Figure~\ref{fig:ch21_nr_frame}; (c) the aggregate of (a) and (b).
Compare with the 20\,Gb/s IMT-2020 requirement.
\item An NR cell at 3.5\,GHz uses the \texttt{DDDSU} pattern with 30\,kHz subcarriers. Compute
the worst-case and average waiting time for an uplink packet that arrives at a random time.
How would a \texttt{DSUUU} pattern change it, and what would it cost?
\item A 5G network is required to deliver 99.999\% reliability within 1\,ms for 32-byte packets.
If one transmission of a mini-slot fails with probability $10^{-3}$ and failures are independent,
how many transmissions are needed? What happens to the latency budget if each retransmission
takes 0.25\,ms, and why might PDCP duplication over two carriers be preferred?
\item List the steps by which a phone switched on in an unfamiliar country finds and joins an
NR standalone cell, from frequency scanning to the first user-plane packet, naming the signals
and channels involved at each step and which of the steps also existed in GSM.
\item \emph{(Simulation)} Extend the power-control simulation in
\lab{book/figscripts/ch21\_figs.py} with an outer loop that adjusts the SIR target to hold a 1\%
block error rate, using a simple BLER-versus-SIR model. Plot the SIR target and the mean
transmit power against speed from 3 to 250\,km/h.
\item \emph{(Simulation)} Generate an LTE downlink signal for a 1.4\,MHz carrier containing the
PSS and SSS of a chosen physical cell identity (TS~36.211 Section~6.11), add a second cell with
a different identity 6\,dB weaker, a frequency offset and noise, and implement the PSS/SSS
search. Measure the probability of detecting both cells against SNR.
\end{enumerate}

\furtherreading
\begin{itemize}
\item E.~Dahlman, S.~Parkvall and J.~Sköld, \emph{5G NR: The Next Generation Wireless Access
Technology}, 2nd ed., Academic Press, 2020; and \emph{4G LTE-Advanced Pro and the Road to 5G},
3rd ed., Academic Press, 2016. Written by Ericsson engineers who took part in the standards;
the clearest explanations of why NR and LTE are designed as they are.
\item S.~Sesia, I.~Toufik and M.~Baker (eds.), \emph{LTE---The UMTS Long Term Evolution: From
Theory to Practice}, 2nd ed., Wiley, 2011. A detailed, physical-layer-heavy reference written by
standards delegates, with good treatment of the reference signals and control channels.
\item H.~Holma and A.~Toskala (eds.), \emph{WCDMA for UMTS: HSPA Evolution and LTE}, 5th ed.,
Wiley, 2010. The standard engineering reference for WCDMA, power control, soft handover and
HSPA, with many field measurements.
\item M.~Mouly and M.-B.~Pautet, \emph{The GSM System for Mobile Communications}, published by the
authors, 1992. The classic description of GSM by two of its designers; still the best book on
the logic of the system.
\item F.~Hillebrand (ed.), \emph{GSM and UMTS: The Creation of Global Mobile Communication},
Wiley, 2002. First-hand accounts by the people who created GSM and UMTS, including the MoU, the
1987 decisions and SMS.
\item A.~J.~Viterbi, \emph{CDMA: Principles of Spread Spectrum Communication}, Addison-Wesley,
1995. The theory behind IS-95 from one of its inventors: power control, soft handoff and
interference-limited capacity.
\item V.~H.~MacDonald, ``The cellular concept,'' \emph{Bell System Technical Journal}, vol.~58,
no.~1, pp.~15--41, 1979. Part of a special issue on AMPS that remains an excellent introduction
to reuse, cell splitting and the AMPS system design.
\item T.~S.~Rappaport, \emph{Wireless Communications: Principles and Practice}, 2nd ed.,
Prentice Hall, 2002. A good textbook treatment of cellular fundamentals and the 1G and 2G
standards.
\item E.~Barkan, E.~Biham and N.~Keller, ``Instant ciphertext-only cryptanalysis of GSM encrypted
communication,'' \emph{Journal of Cryptology}, vol.~21, no.~3, 2008 (first presented at
CRYPTO~2003). How coding-before-ciphering broke A5/2 and undermined A5/1.
\item ITU-R Recommendations M.2083 (IMT-2020 vision, 2015) and M.2410 (minimum requirements,
2017); 3GPP TS 45.002 (GSM multiplexing), TS 25.211--25.214 (WCDMA physical layer), TS 36.211--36.213
(LTE physical layer), TS 38.211--38.214 and 38.306 (NR physical layer and UE capabilities), and
TS 23.501 (5G system architecture). The primary sources, free from the ITU and 3GPP web sites.
\end{itemize}
